% Lesson 14 — Grammar: Polite requests/suggestions: let me + bare infinitive, please/could you (or can you) — Anglais 1ère A — Cours signé OnBuch+
% Programme officiel MINESEC. Compiler avec : tectonic ENA15-grammar-polite-requests-suggestions-let-me-bare-infinit.tex
\newcommand{\DOCMATIERE}{Anglais}
\newcommand{\DOCNIVEAU}{1ère A}
\newcommand{\DOCMODULE}{Chapter 2 : Economic Life and Occupations}
\newcommand{\DOCLECON}{ENA15}
\newcommand{\DOCTITRE}{Lesson 14 — Grammar: Polite requests/suggestions: let me + bare infinitive, please/could you (or can you)}
% OnBuch+ — préambule « cours signés » ENRICHI (compilé avec tectonic / XeLaTeX)
% Système visuel « Pop » d'OnBuch+ : fond crème (jamais blanc), bordures encre
% épaisses, ombres « dures » décalées, titres Archivo Black en capitales.
% Version MATHÉMATIQUES : graphes (pgfplots/tikz), boîtes pédagogiques
% (définition, propriété, méthode, attention, le savais-tu, exercice résolu,
% exercice, TP, bilan), page de garde et signature OnBuch+ ancrée en bas.
%
% Macros attendues dans main.tex AVANT \input{preamble} :
%   \DOCMATIERE  \DOCNIVEAU  \DOCMODULE  \DOCLECON (numéro)  \DOCTITRE
\documentclass[11pt]{article}

\usepackage{fontspec}
\usepackage[english,french]{babel} % français = langue principale ; l'anglais via \eng{..} / textebox
\usepackage[a4paper,margin=2cm,top=3.4cm,bottom=2.8cm,headheight=1.4cm,footskip=1.3cm]{geometry}
\usepackage{amsmath,amssymb,amsfonts}
\usepackage{mathtools} % \xmapsto, \coloneqq... souvent utilisés par les modèles
\usepackage{cancel} % simplifications barrées (bilans de forces)
% \abs{z} (module, valeur absolue) et \norm{u} (norme), utilisés par les modèles
\providecommand{\abs}{}\let\abs\relax\DeclarePairedDelimiter{\abs}{\lvert}{\rvert}
\providecommand{\norm}{}\let\norm\relax\DeclarePairedDelimiter{\norm}{\lVert}{\rVert}
\usepackage[table]{xcolor} % [table] : fournit aussi \rowcolors (alternance de lignes)
\usepackage{pagecolor}
\usepackage{tikz}
\usetikzlibrary{arrows.meta,positioning,calc,shapes.geometric,shapes.misc,shapes.arrows,decorations.pathmorphing,decorations.pathreplacing,decorations.markings,patterns,shadows,mindmap,trees,fit,backgrounds,matrix,intersections,angles,quotes}
\usepackage{pgfplots}
\usepackage[version=4]{mhchem} % équations nucléaires \ce{^{235}_{92}U}
\usepackage[european,siunitx]{circuitikz} % schémas électriques
\pgfplotsset{compat=1.18}
\usepgfplotslibrary{fillbetween,groupplots} % aires entre courbes (calcul intégral)
\usepackage{enumitem}
\usepackage{fancyhdr}
% [most] charge toutes les bibliothèques tcolorbox (listings, minted, raster,
% documentation...) et gonfle inutilement la mémoire TeX sur un document long
% (~300 boîtes) : on ne charge que ce qui est réellement utilisé.
\usepackage[skins,breakable]{tcolorbox}
\usepackage{graphicx}
\usepackage{colortbl}
\usepackage{booktabs}
\usepackage{tabularx}
\usepackage{diagbox} % case d'en-tête diagonale des tableaux à double entrée
% Colonnes extensibles centrée / gauche / droite (C, L, R), souvent utilisées par les modèles
\newcolumntype{C}{>{\centering\arraybackslash}X}
\newcolumntype{L}{>{\raggedright\arraybackslash}X}
\newcolumntype{R}{>{\raggedleft\arraybackslash}X}
\usepackage{array}
\usepackage{multicol}
\usepackage{multirow}
\usepackage{siunitx}
% parse-numbers=false : accepte aussi \SI{2,5 \times 10^{-3}}{...}
\sisetup{locale=FR,output-decimal-marker={,},group-minimum-digits=4,parse-numbers=false}
\DeclareSIUnit\ohm{\ensuremath{\symup{\Omega}}} % U+2126 absent de la police
\DeclareSIUnit\division{div} % réglages d'oscilloscope (ms/div, V/div)
\DeclareSIUnit\tr{tr} % tours (vitesse de rotation, ex. \SI{1500}{\tr\per\minute})
\DeclareSIUnit\molar{M} % concentration molaire (ex. \SI{3}{\molar}, \SI{1}{\micro\molar})
\DeclareSIUnit\an{an} % année (le modèle confond parfois avec \year, un compteur TeX) — ex. \SI{5}{\kilo\meter\per\an}
\DeclareSIUnit\FCFA{FCFA} % franc CFA (ex. \SI{500}{\FCFA\per\kilo\gram})
\DeclareSIUnit\degreeBrix{\textdegree Brix} % teneur en sucre d'un sirop/jus (réfractométrie)
\DeclareSIUnit\month{mois} % le modèle utilise parfois \month, un compteur TeX, comme unité de durée
\usepackage{qrcode}
% \degree : commande fréquemment utilisée par les modèles pour un angle ou une
% température en dehors de \SI{}{} (non fournie par les paquets chargés).
\providecommand{\degree}{^{\circ}}
% \qed (fin de démonstration), utilisé par les modèles sans amsthm
\providecommand{\qed}{\hfill$\square$}
% \barre : double barre des valeurs interdites dans un tableau de variations (tkz-tab)
\providecommand{\barre}{\ensuremath{\|}}
% environnement proof (amsthm non chargé) utilisé par les modèles
\ifdefined\proof\else\newenvironment{proof}[1][Démonstration]{\par\noindent\textit{#1.}\ }{\qed\par}\fi
\usepackage{lastpage}
\usepackage{hyperref}
\hypersetup{hidelinks}

% ============================================================
% Palette « Pop » OnBuch+ — mêmes hex que la classe P (plus_ui.dart)
% ============================================================
\definecolor{poporange}{HTML}{F59321}
\definecolor{popdark}{HTML}{E07A0C}
\definecolor{popcream}{HTML}{FFF6E8}
\definecolor{popink}{HTML}{1C1714}
\definecolor{poppurple}{HTML}{7A5AE0}
\definecolor{popgreen}{HTML}{2E9E6A}
\definecolor{poppink}{HTML}{E0457B}
\definecolor{popblue}{HTML}{2F7DE1}
\definecolor{popgold}{HTML}{C98A12}
\definecolor{popmuted}{HTML}{8A7F76}
\definecolor{popline}{HTML}{EADFD2}
\definecolor{popgray}{HTML}{8A7F76}
\colorlet{popgrey}{popgray}
% Alias tolérants (le modèle nomme parfois les couleurs différemment)
\colorlet{popwhite}{white}
\colorlet{popblack}{popink}
\colorlet{popred}{poppink}
\colorlet{popyellow}{popgold}
% Teintes claires (fonds de tableaux/encadrés) — le modèle nomme parfois
% les couleurs avec un suffixe L (ex. popblueL) sans les avoir définies.
\colorlet{poporangeL}{poporange!20}
\colorlet{popdarkL}{popdark!20}
\colorlet{poppurpleL}{poppurple!20}
\colorlet{popgreenL}{popgreen!20}
\colorlet{poppinkL}{poppink!20}
\colorlet{popblueL}{popblue!20}
\colorlet{popgoldL}{popgold!20}
\colorlet{popredL}{popred!20}
\colorlet{popgrayL}{popgray!20}
% Teintes claires pour fonds de tableaux / zones
\colorlet{popblueL}{popblue!10!white}
\colorlet{poporangeL}{poporange!14!white}
\colorlet{popgreenL}{popgreen!12!white}
\colorlet{poppurpleL}{poppurple!12!white}
\colorlet{poppinkL}{poppink!12!white}
\colorlet{popgoldL}{popgold!14!white}

\pagecolor{popcream}

% ============================================================
% Typographie
% ============================================================
\defaultfontfeatures{Ligatures=TeX}
\setmainfont{PlusJakartaSans-Medium.ttf}[Path=../../../fonts/,BoldFont=PlusJakartaSans-Bold.ttf]
% Symboles, API et emojis absents de la police principale : repli sur DejaVu Sans (caractères actifs)
\newfontface\symfont{DejaVuSans.ttf}[Path=../../../fonts/]
\makeatletter
\newcommand\symactive[1]{\catcode#1=13 \begingroup\lccode`\~=#1 \lowercase{\endgroup\def~}{{\symfont\char#1}}}
\newcommand\symblank[1]{\catcode#1=13 \begingroup\lccode`\~=#1 \lowercase{\endgroup\def~}{}}
\newcommand\symhyph[1]{\catcode#1=13 \begingroup\lccode`\~=#1 \lowercase{\endgroup\def~}{\mbox{-}}}
\newcount\symc
\newcommand\symrange[2]{\symc=#1 \@whilenum\symc<#2 \do{\expandafter\symactive\expandafter{\the\symc}\advance\symc by 1 }}
\newcommand\symblankrange[2]{\symc=#1 \@whilenum\symc<#2 \do{\expandafter\symblank\expandafter{\the\symc}\advance\symc by 1 }}
\makeatother
\symrange{"0250}{"02FF}\symactive{"2713}\symactive{"2714}\symactive{"2717}\symactive{"2718}\symactive{"2610}\symactive{"2611}\symactive{"2612}
\symhyph{"2011}\symhyph{"2E1A}\symblankrange{"1F300}{"1FAFF}\symblank{"FE0F}
% Maths en sans-serif (Fira Math) avec chiffres et lettres droites de Plus
% Jakarta, pour que les formules se fondent dans le texte.
\usepackage[math-style=ISO,bold-style=ISO]{unicode-math}
\setmathfont{FiraMath-Regular.otf}[Path=../../../fonts/]
\setmathfont{PlusJakartaSans-Medium.ttf}[Path=../../../fonts/,range={up/num,up/{Latin,latin},\mathup}]
\usepackage{icomma} % virgule décimale sans espace en mode math
% Caractères mathématiques Unicode absents des polices de texte (Plus Jakarta,
% Archivo Black) : rendus par leur équivalent mathématique, en texte comme en
% formule — sinon ils s'affichent en carré vide (ex. « Arithmétique dans ℤ »).
\usepackage{newunicodechar}
\newunicodechar{ℝ}{\ensuremath{\mathbb{R}}}
\newunicodechar{ℤ}{\ensuremath{\mathbb{Z}}}
\newunicodechar{ℂ}{\ensuremath{\mathbb{C}}}
\newunicodechar{ℕ}{\ensuremath{\mathbb{N}}}
\newunicodechar{ℚ}{\ensuremath{\mathbb{Q}}}
\newunicodechar{𝔻}{\ensuremath{\mathbb{D}}}
\newunicodechar{α}{\ensuremath{\alpha}}
\newunicodechar{β}{\ensuremath{\beta}}
\newunicodechar{γ}{\ensuremath{\gamma}}
\newunicodechar{δ}{\ensuremath{\delta}}
\newunicodechar{ε}{\ensuremath{\varepsilon}}
\newunicodechar{θ}{\ensuremath{\theta}}
\newunicodechar{λ}{\ensuremath{\lambda}}
\newunicodechar{μ}{\ensuremath{\mu}}
\newunicodechar{σ}{\ensuremath{\sigma}}
\newunicodechar{τ}{\ensuremath{\tau}}
\newunicodechar{φ}{\ensuremath{\varphi}}
\newunicodechar{ψ}{\ensuremath{\psi}}
\newunicodechar{ω}{\ensuremath{\omega}}
\newunicodechar{Σ}{\ensuremath{\Sigma}}
% majuscules produites par \MakeUppercase dans les titres (α-aminés -> Α)
\newunicodechar{Α}{\ensuremath{\alpha}}
\newunicodechar{Β}{\ensuremath{\beta}}
\newunicodechar{Γ}{\ensuremath{\Gamma}}
\newunicodechar{Δ}{\ensuremath{\Delta}}
\newunicodechar{Ε}{\ensuremath{\varepsilon}}
\newunicodechar{Θ}{\ensuremath{\Theta}}
\newunicodechar{Λ}{\ensuremath{\Lambda}}
\newunicodechar{Μ}{\ensuremath{\mu}}
\newunicodechar{Τ}{\ensuremath{\tau}}
\newunicodechar{Φ}{\ensuremath{\Phi}}
\newunicodechar{Ψ}{\ensuremath{\Psi}}
\newunicodechar{Ω}{\ensuremath{\Omega}}
\newunicodechar{↦}{\ensuremath{\mapsto}}
\newunicodechar{∈}{\ensuremath{\in}}
\newunicodechar{∉}{\ensuremath{\notin}}
\newunicodechar{∧}{\ensuremath{\wedge}}
\newunicodechar{⇔}{\ensuremath{\Leftrightarrow}}
\newunicodechar{⟺}{\ensuremath{\Longleftrightarrow}}
\newunicodechar{⇒}{\ensuremath{\Rightarrow}}
\newunicodechar{⟹}{\ensuremath{\Longrightarrow}}
\newunicodechar{∘}{\ensuremath{\circ}}
\newunicodechar{∪}{\ensuremath{\cup}}
\newunicodechar{∩}{\ensuremath{\cap}}
\newunicodechar{≡}{\ensuremath{\equiv}}
\newunicodechar{⊂}{\ensuremath{\subset}}
\newunicodechar{⊥}{\ensuremath{\perp}}
\newunicodechar{∃}{\ensuremath{\exists}}
\newunicodechar{∀}{\ensuremath{\forall}}
\newunicodechar{∛}{\ensuremath{\sqrt[3]{\,}}}
\newunicodechar{∜}{\ensuremath{\sqrt[4]{\,}}}
\newunicodechar{⃗}{\ensuremath{{}^{\to}}}
\newunicodechar{ⁿ}{\ifmmode^{n}\else\textsuperscript{\ensuremath{n}}\fi}
\newunicodechar{ˣ}{\ifmmode^{x}\else\textsuperscript{\ensuremath{x}}\fi}
\newunicodechar{ᵇ}{\ifmmode^{b}\else\textsuperscript{\ensuremath{b}}\fi}
\newunicodechar{ʳ}{\ifmmode^{r}\else\textsuperscript{\ensuremath{r}}\fi}
\newunicodechar{ᵘ}{\ifmmode^{u}\else\textsuperscript{\ensuremath{u}}\fi}
\newunicodechar{ʸ}{\ifmmode^{y}\else\textsuperscript{\ensuremath{y}}\fi}
\newunicodechar{ᵃ}{\ifmmode^{a}\else\textsuperscript{\ensuremath{a}}\fi}
\newunicodechar{ᵉ}{\ifmmode^{e}\else\textsuperscript{\ensuremath{e}}\fi}
\newunicodechar{ᵏ}{\ifmmode^{k}\else\textsuperscript{\ensuremath{k}}\fi}
\newunicodechar{ᵖ}{\ifmmode^{p}\else\textsuperscript{\ensuremath{p}}\fi}
\newunicodechar{ᴺ}{\ifmmode^{N}\else\textsuperscript{\ensuremath{N}}\fi}
\newunicodechar{ᶜ}{\ifmmode^{c}\else\textsuperscript{\ensuremath{c}}\fi}
\newunicodechar{ʰ}{\ifmmode^{h}\else\textsuperscript{\ensuremath{h}}\fi}
\newunicodechar{ᵠ}{\ifmmode^{q}\else\textsuperscript{\ensuremath{q}}\fi}
\newunicodechar{ₙ}{\ifmmode_{n}\else\textsubscript{\ensuremath{n}}\fi}
\newunicodechar{ₐ}{\ifmmode_{a}\else\textsubscript{\ensuremath{a}}\fi}
\newunicodechar{ᵢ}{\ifmmode_{i}\else\textsubscript{\ensuremath{i}}\fi}
\newunicodechar{ᵣ}{\ifmmode_{r}\else\textsubscript{\ensuremath{r}}\fi}
\newunicodechar{ₖ}{\ifmmode_{k}\else\textsubscript{\ensuremath{k}}\fi}
\newunicodechar{ᵦ}{\ifmmode_{\beta}\else\textsubscript{\ensuremath{\beta}}\fi}
\newunicodechar{ₓ}{\ifmmode_{x}\else\textsubscript{\ensuremath{x}}\fi}
\newfontfamily\popdisplay{ArchivoBlack-Regular.ttf}[Path=../../../fonts/]
\newfontfamily\poplogo{SpaceGrotesk-Bold.ttf}[Path=../../../fonts/]
\newfontfamily\popsemi{PlusJakartaSans-SemiBold.ttf}[Path=../../../fonts/]
\newfontfamily\popxbold{PlusJakartaSans-ExtraBold.ttf}[Path=../../../fonts/]
\newfontfamily\popxboldls{PlusJakartaSans-ExtraBold.ttf}[Path=../../../fonts/,LetterSpace=3.0]

\setlist[enumerate]{leftmargin=1.7em,itemsep=5pt,topsep=4pt}
\setlist[itemize]{leftmargin=1.7em,itemsep=4pt,topsep=4pt,label={\color{poporange}\textbullet}}
\setlength{\parindent}{0pt}
\setlength{\parskip}{5pt}
\renewcommand{\arraystretch}{1.25}

% pgfplots : style Pop par défaut
\pgfplotsset{
  every axis/.append style={axis line style={popink,line width=1pt},tick style={popink},
    grid style={popline,line width=0.5pt},label style={font=\small},tick label style={font=\footnotesize}},
  popcurve/.style={poporange,line width=1.8pt,no markers,smooth},
}
% Même style disponible dans un \draw TikZ simple (hors pgfplots)
\tikzset{popcurve/.style={poporange,line width=1.8pt}}
\tikzset{popgrid/.style={popline,very thin}}
\colorlet{popcurve}{poporange} % le nom du style est parfois utilisé comme couleur

% ============================================================
% Pill / squiggle
% ============================================================
\newcommand{\poppill}[3]{%
  \tikz[baseline=-0.55ex]\node[draw=popink,line width=1.2pt,rounded corners=2.6pt,%
    fill=#2,inner xsep=6.5pt,inner ysep=3pt]{%
    \popxboldls\fontsize{7}{7}\selectfont\color{#3}\MakeUppercase{#1}};%
}
\newcommand{\popsquiggle}[1]{%
  \tikz[baseline=-0.4ex]{
    \draw[#1,line width=2.6pt,line cap=round]
      plot [smooth] coordinates {(0,0.09) (0.34,0.01) (0.68,0.21) (1.02,0.01) (1.36,0.10)};
  }%
}
% Mot-clé surligné (marqueur orange sous le texte)
\newcommand{\cle}[1]{{\popxbold\color{popdark}#1}}

% ============================================================
% En-tête / pied
% ============================================================
\pagestyle{fancy}
\fancyhf{}
\renewcommand{\headrulewidth}{0pt}
\renewcommand{\footrulewidth}{0pt}
\lhead{\raisebox{-0.15ex}{\poplogo\fontsize{13}{13}\selectfont\color{popink}OnBuch\color{poporange}+}}
\rhead{\poppill{\DOCMATIERE\ \textbullet\ \DOCNIVEAU\ \textbullet\ Leçon \DOCLECON}{white}{popink}}
\lfoot{\popxbold\fontsize{7.6}{7.6}\selectfont\color{popmuted}\MakeUppercase{Cours signé OnBuch+}\ \textbullet\ {\popsemi\DOCTITRE}}
\rfoot{\tikz[baseline=-0.6ex]\node[rounded corners=8.5pt,fill=popink,minimum height=17pt,minimum width=17pt,inner xsep=5pt,inner sep=0]{\popxbold\fontsize{8}{8}\selectfont\color{popcream}\thepage\,/\,\pageref*{LastPage}};}

% ============================================================
% Titres
% ============================================================
\newcommand{\chaptitre}[2]{%
  \begin{tcolorbox}[enhanced,colback=poporange,colframe=popink,boxrule=2.6pt,arc=11pt,
      drop shadow=popink,left=15pt,right=15pt,top=12pt,bottom=11pt,before skip=3pt,after skip=18pt]
    \poppill{#2}{popink}{popcream}\par\vspace{9pt}
    {\raggedright\hyphenpenalty=10000\exhyphenpenalty=10000\popdisplay\fontsize{22}{24}\selectfont\color{popcream}\MakeUppercase{#1}\par}
  \end{tcolorbox}
}
% Section numérotée : pastille orange avec le numéro + titre Archivo
\newcounter{popsec}
\newcommand{\coursec}[1]{%
  \stepcounter{popsec}\setcounter{popsub}{0}%
  \par\vspace{16pt}\noindent
  \begin{minipage}{\linewidth}
  \tikz[baseline=-0.7ex]\node[draw=popink,line width=1.6pt,fill=poporange,rounded corners=4pt,minimum size=22pt,inner sep=0]{\popdisplay\fontsize{12}{12}\selectfont\color{popcream}\thepopsec};%
  \hspace{8pt}{\popdisplay\fontsize{15}{17}\selectfont\color{popink}\MakeUppercase{#1}}\par
  \vspace{4pt}\noindent{\color{poporange}\rule{36pt}{3.6pt}}
  \end{minipage}\par\nopagebreak\vspace{8pt}
}
\newcounter{popsub}
\newcommand{\courssub}[1]{%
  \stepcounter{popsub}%
  \par\vspace{9pt}\noindent{\popxbold\fontsize{11.5}{13}\selectfont\color{popink}{\color{poporange}\thepopsec.\thepopsub}\hspace{6pt}#1}\par\nopagebreak\vspace{4pt}
}

% ============================================================
% Boîtes pédagogiques — même recette : fond blanc, bordure épaisse colorée,
% ombre dure encre, badge plein. Titre optionnel [#1].
% ============================================================
\tcbset{popbase/.style={breakable,enhanced,colback=white,boxrule=2.3pt,arc=9pt,
  shadow={3pt}{-3pt}{0pt}{popink},left=14pt,right=14pt,top=11pt,bottom=11pt,before skip=11pt,after skip=15pt,
  fontupper=\popsemi\fontsize{10.6}{14.5}\selectfont\color{popink},
  fontlower=\popsemi\fontsize{10.6}{14.5}\selectfont\color{popink}}}
% Coche dessinée (le glyphe ✓ n'existe pas dans les polices du document)
\newcommand{\popcheck}[1]{\tikz[baseline=-0.5ex]\draw[#1,line width=1.6pt,line cap=round,line join=round](-0.13,0.0)--(-0.04,-0.09)--(0.13,0.1);}
\providecommand{\coche}{\popcheck{popgreen}}
\providecommand{\resultat}[1]{\par\smallskip\noindent\fcolorbox{popgreen}{popgreenL}{\parbox{\dimexpr\linewidth-2\fboxsep-2\fboxrule\relax}{\textbf{Résultat.} #1}}\par\smallskip}
\newcommand{\popboxhead}[3]{\poppill{#1}{#2}{white}\ifx\relax#3\relax\else\hspace{6pt}{\popxbold\fontsize{10.6}{12}\selectfont\color{popink}#3}\fi\par\vspace{7pt}}

\newtcolorbox{aretenir}[1][]{popbase,colframe=popgreen,before upper={\popboxhead{À retenir}{popgreen}{#1}}}
% Texte en anglais (document de lecture, dialogue, extrait) : cadre
% d'étude. Un titre optionnel (source, auteur) s'affiche dans l'en-tête.
\newtcolorbox{textebox}[1][]{popbase,colframe=popblue,before upper={\popboxhead{Text}{popblue}{#1}\begin{otherlanguage}{english}},after upper={\end{otherlanguage}}}
% Marques ✓ / ✗ (forme correcte / fautive) souvent utilisées par les modèles
\providecommand{\cross}{\textcolor{poppink}{\ensuremath{\times}}}
\providecommand{\xmark}{\cross}\providecommand{\crossmark}{\cross}\providecommand{\cmark}{\ensuremath{\checkmark}}
% Ligne de réponse / justification à compléter (inventée par les modèles)
\providecommand{\justification}[1]{\par\noindent\textit{Justification~:} \dotfill\par}
% \textipa (package tipa non chargé) : la police ne couvre pas l'API -> simple passage
\providecommand{\textipa}[1]{#1}
% Soulignement (ulem non chargé) : \uline, \ul
\providecommand{\uline}[1]{\underline{#1}}\providecommand{\ul}[1]{\underline{#1}}
% \glossaire{\item ...} inventée par les modèles : liste de glossaire
\providecommand{\glossaire}[1]{\begin{itemize}#1\end{itemize}}
% \alert (beamer) inventée par les modèles : texte mis en évidence
\providecommand{\alert}[1]{\textbf{\textcolor{poppink}{#1}}}
% variantes ASCII d'ordinaux écrites en macro
\providecommand{\iere}{\textsuperscript{re}}\providecommand{\ieme}{\textsuperscript{e}}\providecommand{\ere}{\textsuperscript{re}}\providecommand{\eme}{\textsuperscript{e}}\providecommand{\er}{\textsuperscript{er}}\providecommand{\ie}{\textsuperscript{e}}
% Ordinaux français écrits en macro par les modèles : 1\ière, 3\ième
\newcommand{\ière}{\textsuperscript{re}}\newcommand{\ième}{\textsuperscript{e}}
% \exemple (inventée par les modèles) : étiquette « Exemple : »
\providecommand{\exemple}{\textbf{Exemple~:}\ }
% \square utilisable aussi hors mode math (cases à cocher)
\AtBeginDocument{\let\squareMath\square\renewcommand{\square}{\ensuremath{\squareMath}}}
% Mot ou phrase en anglais à l'intérieur d'un paragraphe français
\newcommand{\eng}[1]{\ifmmode\text{\textit{#1}}\else\begingroup\selectlanguage{english}\itshape #1\endgroup\fi}
\let\esp\eng % alias toléré
% flèches d'intonation inventées par les modèles
\providecommand{\textsearrow}{}\renewcommand{\textsearrow}{\ensuremath{\searrow}}\providecommand{\textnearrow}{}\renewcommand{\textnearrow}{\ensuremath{\nearrow}}
% \fr{..} : mot français (inventé par les modèles)
\providecommand{\fr}[1]{\textit{#1}}
\newtcolorbox{exemplebox}[1][]{popbase,colframe=poporange,before upper={\popboxhead{Exemple}{poporange}{#1}}}
\newtcolorbox{definition}[1][]{popbase,colframe=popblue,before upper={\popboxhead{Définition}{popblue}{#1}}}
\newtcolorbox{propriete}[1][]{popbase,colframe=poppurple,before upper={\popboxhead{Propriété}{poppurple}{#1}}}
\newtcolorbox{methode}[1][]{popbase,colframe=popgold,before upper={\popboxhead{Méthode}{popgold}{#1}}}
\newtcolorbox{attention}[1][]{popbase,colframe=poppink,before upper={\popboxhead{Attention}{poppink}{#1}}}
\newtcolorbox{experience}[1][]{popbase,colframe=popblue,colback=popblueL,before upper={\popboxhead{Expérience}{popblue}{#1}}}
% « Le savais-tu ? » : carte violette pleine (contexte, culture, Cameroun)
\newtcolorbox{savaistu}[1][]{popbase,colback=poppurple,colframe=popink,
  fontupper=\popsemi\fontsize{10.4}{14.2}\selectfont\color{white},
  before upper={\poppill{Le savais-tu\,?}{white}{poppurple}\ifx\relax#1\relax\else\hspace{6pt}{\popxbold\color{white}#1}\fi\par\vspace{7pt}}}
% Exercice résolu : énoncé en haut, \tcblower puis la solution sur fond crème
\newcounter{popexo}\newcounter{popexr}
\newtcolorbox{exoresolu}[1][]{popbase,colframe=poporange,colbacklower=popcream,
  segmentation style={popink,line width=1.2pt,dashed},
  before upper={\stepcounter{popexr}\popboxhead{Exercice résolu \thepopexr}{poporange}{#1}},
  before lower={\poppill{Solution}{popink}{popcream}\par\vspace{6pt}}}
% Exercice (non résolu en place) — la correction est dans la partie Corrigés
\newtcolorbox{exercice}[1][]{popbase,colframe=popink,
  before upper={\stepcounter{popexo}\popboxhead{Exercice \thepopexo}{popink}{#1}}}
% Corrigé
\newtcolorbox{corrige}[1][]{popbase,colframe=popgreen,colback=popgreenL,
  before upper={\popboxhead{Corrigé}{popgreen}{#1}}}
% Objectifs / prérequis / bilan
\newtcolorbox{objectifs}{popbase,colframe=popink,colback=poporangeL,
  before upper={\popboxhead{Objectifs de la leçon}{popink}{}}}
\newtcolorbox{prerequis}{popbase,colframe=popink,colback=popblueL,
  before upper={\popboxhead{Prérequis}{popblue}{}}}
\newtcolorbox{bilan}{popbase,colframe=popink,colback=white,boxrule=2.8pt,
  before upper={\popboxhead{Fiche bilan}{poporange}{L'essentiel à connaître pour l'examen}}}
% Figure encadrée avec légende
\newtcolorbox{popfigure}[1][]{enhanced,breakable=false,colback=white,colframe=popline,boxrule=1.4pt,arc=8pt,
  left=10pt,right=10pt,top=10pt,bottom=8pt,before skip=10pt,after skip=12pt,halign=center,
  lower separated=false,halign lower=center,
  fontlower=\popsemi\fontsize{9}{11.5}\selectfont\color{popmuted},
  #1}
\newcommand{\legende}[1]{\par\vspace{6pt}{\popsemi\fontsize{9}{11.5}\selectfont\color{popmuted}#1}}

% ============================================================
% Signature OnBuch+
% ============================================================
\newcommand{\poplogoinline}[1]{{\poplogo\fontsize{#1}{#1}\selectfont\color{popink}OnBuch\color{poporange}+}}
\newcommand{\signatureonbuch}{%
  \begin{tcolorbox}[enhanced,colback=popink,colframe=popink,boxrule=2.2pt,arc=10pt,
      drop shadow=poporange,left=14pt,right=14pt,top=10pt,bottom=10pt,before skip=0pt,after skip=0pt]
    \tikz[baseline=-0.7ex]\node[circle,fill=popgreen,draw=popcream,line width=1.3pt,minimum size=22pt,inner sep=0]{\popcheck{white}};%
    \hspace{9pt}\begin{minipage}[c]{0.62\linewidth}
      {\popxbold\fontsize{12}{13}\selectfont\color{popcream}Signé\ \textbullet\ {\poplogo OnBuch\color{poporange}+}}\par\vspace{2pt}
      {\raggedright\popsemi\fontsize{8}{10.5}\selectfont\color{popline}\MakeUppercase{Équipe pédagogique OnBuch+}\\\MakeUppercase{Conforme au programme officiel MINESEC}\par}
    \end{minipage}\hfill
    {\poplogo\fontsize{22}{22}\selectfont\color{popcream}OnBuch\color{poporange}+}
  \end{tcolorbox}%
}

% ============================================================
% Page de garde
% #1 titre · #2 matière · #3 niveau · #4 module · #5 n° leçon · #6 durée ·
% #7 description / objectifs (texte ou itemize)
% ============================================================
\newcommand{\pagegarde}[7]{%
  \thispagestyle{empty}
  \newgeometry{a4paper,margin=1.7cm,top=1.6cm,bottom=1.4cm}%
  \begin{tikzpicture}[remember picture,overlay]
    \fill[poporange!18] ([xshift=-3.2cm,yshift=-3.6cm]current page.north east) circle (4.6cm);
    \fill[poppurple!14] ([xshift=2.4cm,yshift=7.6cm]current page.south west) circle (3.8cm);
  \end{tikzpicture}%
  {\poplogo\fontsize{30}{30}\selectfont\color{popink}OnBuch\color{poporange}+}\hfill
  \raisebox{0.6ex}{\poppill{Cours signé \textbullet\ Premium}{popink}{popcream}}\par
  \vspace{1pt}\popsquiggle{poppurple}\par
  \vspace{8pt}
  \begin{tcolorbox}[enhanced,colback=poporange,colframe=popink,boxrule=2.8pt,arc=13pt,
      drop shadow=popink,left=18pt,right=18pt,top=12pt,bottom=13pt,after skip=10pt]
    \poppill{#2 \textbullet\ #3}{popink}{popcream}\hspace{5pt}\poppill{#4}{white}{popink}\par\vspace{8pt}
    {\popdisplay\fontsize{14}{14}\selectfont\color{popink}LEÇON #5}\par\vspace{4pt}
    {\raggedright\hyphenpenalty=10000\exhyphenpenalty=10000\popdisplay\fontsize{25}{27}\selectfont\color{popcream}\MakeUppercase{#1}\par}
    \vspace{8pt}
    {\popxbold\fontsize{9.5}{9.5}\selectfont\color{popink}\MakeUppercase{Durée conseillée : #6}}
  \end{tcolorbox}
  \begin{tcolorbox}[enhanced,colback=white,colframe=popink,boxrule=2.2pt,arc=10pt,
      drop shadow=popink,left=16pt,right=16pt,top=10pt,bottom=10pt,after skip=8pt]
    \poppill{Dans cette leçon}{poppurple}{white}\par\vspace{5pt}
    \popsemi\fontsize{9.3}{12.6}\selectfont\color{popink}#7
  \end{tcolorbox}
  \noindent\begin{minipage}[t]{0.487\linewidth}
    \begin{tcolorbox}[enhanced,colback=popgreen,colframe=popink,boxrule=2.2pt,arc=10pt,
        drop shadow=popink,left=11pt,right=11pt,top=10pt,bottom=10pt,height=3.1cm,valign=center]
      \tcbox[colback=white,colframe=popink,boxrule=1.3pt,arc=5pt,boxsep=0pt,nobeforeafter,
        left=3pt,right=3pt,top=3pt,bottom=3pt,valign=center]{\qrcode[height=1.9cm]{https://play.google.com/store/apps/details?id=com.luvvix.onbuch&hl=fr}}%
      \hspace{8pt}\begin{minipage}[c]{0.5\linewidth}
        {\popdisplay\fontsize{11}{12.5}\selectfont\color{white}\MakeUppercase{Télécharge OnBuch}}\par\vspace{4pt}
        {\raggedright\popsemi\fontsize{8.4}{11}\selectfont\color{white}Gratuit \textbullet\ Google Play\\Tuteur IA, annales, concours, résultats\par}
      \end{minipage}
    \end{tcolorbox}
  \end{minipage}\hfill
  \begin{minipage}[t]{0.487\linewidth}
    \begin{tcolorbox}[enhanced,colback=poppurple,colframe=popink,boxrule=2.2pt,arc=10pt,
        drop shadow=popink,left=11pt,right=11pt,top=10pt,bottom=10pt,height=3.1cm,valign=center]
      \tikz[baseline=-0.7ex]\node[circle,fill=white,draw=popink,line width=1.3pt,minimum size=1.6cm,inner sep=0]{\popdisplay\fontsize{24}{24}\selectfont\color{poppurple}+};%
      \hspace{8pt}\begin{minipage}[c]{0.6\linewidth}
        {\popdisplay\fontsize{11}{12.5}\selectfont\color{white}\MakeUppercase{Passe à OnBuch+}}\par\vspace{4pt}
        {\raggedright\popsemi\fontsize{8.4}{11}\selectfont\color{white}Tous les cours signés, Tuteur IA illimité, fiches PDF — onglet OnBuch+ de l'app\par}
      \end{minipage}
    \end{tcolorbox}
  \end{minipage}
  \vspace{8pt}
  \popcontenu
  \vfill
  \signatureonbuch
  \restoregeometry
  \newpage
}

% Rangée « Ce cours contient » (page de garde)
\newcommand{\poptile}[3]{% #1 couleur #2 gros texte #3 légende
  \begin{tcolorbox}[enhanced,colback=#1,colframe=popink,boxrule=2pt,arc=9pt,shadow={2.5pt}{-2.5pt}{0pt}{popink},
      left=6pt,right=6pt,top=9pt,bottom=9pt,halign=center,valign=center,height=1.75cm,width=\linewidth,nobeforeafter]
    {\popdisplay\fontsize{12.5}{13}\selectfont\color{white}\MakeUppercase{#2}}\par\vspace{3pt}
    {\centering\popsemi\fontsize{7.8}{9.5}\selectfont\color{white}#3\par}
  \end{tcolorbox}}
\newcommand{\popcontenu}{%
  \par\noindent{\popxboldls\fontsize{7.5}{7.5}\selectfont\color{popmuted}CE COURS SIGNÉ CONTIENT}\par\vspace{7pt}
  \noindent\begin{minipage}[t]{0.235\linewidth}\poptile{popblue}{Cours}{complet, illustré, pas à pas}\end{minipage}\hfill
  \begin{minipage}[t]{0.235\linewidth}\poptile{poporange}{Méthodes}{et exercices résolus}\end{minipage}\hfill
  \begin{minipage}[t]{0.235\linewidth}\poptile{poppink}{Exercices}{type examen + corrigés}\end{minipage}\hfill
  \begin{minipage}[t]{0.235\linewidth}\poptile{popgreen}{Fiche bilan}{l'essentiel à retenir}\end{minipage}\par}

% Sommaire
\newtcolorbox{sommaire}{enhanced,colback=white,colframe=popink,boxrule=2.2pt,arc=10pt,
  drop shadow=popink,left=18pt,right=18pt,top=13pt,bottom=13pt,before skip=6pt,after skip=20pt,
  before upper={\poppill{Sommaire}{poppurple}{white}\par\vspace{9pt}\popsemi\fontsize{10.6}{14.5}\selectfont\color{popink}}}

% Signature de fin de cours — poussée tout en bas de la dernière page
\newcommand{\findecours}{%
  \par\vfill\nopagebreak
  \begin{center}\popsquiggle{poporange}\end{center}\vspace{4pt}
  \signatureonbuch
}
\AtBeginDocument{\def\llbracket{\mathopen{[\mkern-3.5mu[}}\def\rrbracket{\mathclose{]\mkern-3.5mu]}}} % intervalles d'entiers (glyphe ⟦ absent de la police)
% Glyphes absents de Fira Math : redéfinis à partir de glyphes disponibles
\AtBeginDocument{%
  \def\setminus{\mathbin{\backslash}}\let\smallsetminus\setminus
  \def\vdots{\mathord{\vbox{\baselineskip3.2pt\lineskiplimit0pt\kern4pt\hbox{.}\hbox{.}\hbox{.}}}}%
  \def\ddots{\mathinner{\mkern1mu\raise6.5pt\hbox{.}\mkern2mu\raise3.6pt\hbox{.}\mkern2mu\raise0.7pt\hbox{.}\mkern1mu}}%
  \def\triangle{\mathord{\scalebox{0.9}{$\symup{\Delta}$}}}%
  \def\nabla{\mathord{\raisebox{\depth}{\scalebox{1}[-1]{$\symup{\Delta}$}}}}%
  \def\checkmark{\popcheck{popdark}}%
  \def\star{\mathbin{\ast}}\def\bigstar{\mathord{\ast}}%
  \def\diamond{\mathbin{\scalebox{0.6}{$\Diamond$}}}%
  \def\bigcup{\mathop{\mathchoice{\raisebox{-0.3em}{\scalebox{1.7}{$\cup$}}}{\scalebox{1.25}{$\cup$}}{\cup}{\cup}}\displaylimits}%
  \def\bigcap{\mathop{\mathchoice{\raisebox{-0.3em}{\scalebox{1.7}{$\cap$}}}{\scalebox{1.25}{$\cap$}}{\cap}{\cap}}\displaylimits}%
  \def\lesseqqgtr{\mathrel{\lessgtr}}\def\bigoplus{\mathop{\oplus}\displaylimits}%
}
\providecommand{\ssi}{si et seulement si} % abréviation fréquente
\AtBeginDocument{\providecommand{\wideparen}{\overparen}} % arc de cercle

%% FIGFIX-BEGIN v5 — correctifs globaux des figures (docs/onbuch-generation/tools/figfix)
\usepackage{adjustbox}\usepackage{etoolbox}\tcbuselibrary{raster}
% toute figure TikZ/pgfplots plus large que la ligne est réduite à la largeur disponible
\BeforeBeginEnvironment{tikzpicture}{\begin{adjustbox}{max width=\linewidth}}
\AfterEndEnvironment{tikzpicture}{\end{adjustbox}}
% légendes pgfplots lisibles par-dessus les courbes
\pgfplotsset{every axis/.append style={legend style={fill=white,fill opacity=0.92,text opacity=1,draw=black!15,inner sep=3pt}}}
% étiquettes d'arêtes (syntaxe edge["..."]) avec fond blanc
\tikzset{every edge quotes/.append style={fill=white,inner sep=1.2pt}}
%% FIGFIX-END
\begin{document}

\pagegarde{Lesson 14 — Grammar: Polite requests/suggestions: let me + bare infinitive, please/could you (or can you)}{Anglais}{1ère A}{Chapter 2 : Economic Life and Occupations}{ENA15}{2 h}{Cette leçon de grammaire présente les structures de demandes polies et de suggestions en anglais : let me + base verbale, please, could you et can you. À travers l'observation d'exemples authentiques, des règles claires et des exercices progressifs, les élèves apprennent à formuler des requêtes courtoises adaptées à la vie économique et professionnelle.
\vspace{4pt}
\begin{multicols}{2}\small
\begin{itemize}
\item À la fin de cette leçon, je sais identifier la structure let me + base verbale dans un énoncé
\item À la fin de cette leçon, je sais construire une suggestion polie avec let me + infinitif sans to
\item À la fin de cette leçon, je sais formuler une demande polie avec could you et can you + base verbale
\item À la fin de cette leçon, je sais placer correctement please dans une demande polie
\item À la fin de cette leçon, je sais distinguer les niveaux de politesse (can you < could you < could you ... please)
\item À la fin de cette leçon, je sais éviter les erreurs fréquentes des francophones (to après let, traduction mot à mot de 'pouvez-vous')
\end{itemize}\end{multicols}}

\begin{sommaire}
\begin{multicols}{2}
\begin{enumerate}
\item Observation et découverte
\item Règle 1 : Let me + base verbale (suggestions et offres d'aide)
\item Règle 2 : Could you / Can you + base verbale (demandes polies)
\item Comparaison français/anglais et pièges des francophones
\item Entraînement guidé et usage en contexte économique
\item Activité d'intégration
\item Méthodes et exercices résolus
\item Exercices
\item Corrigés des exercices
\item Fiche bilan
\end{enumerate}
\end{multicols}
\end{sommaire}

\begin{objectifs}
\begin{itemize}
\item À la fin de cette leçon, je sais identifier la structure let me + base verbale dans un énoncé
\item À la fin de cette leçon, je sais construire une suggestion polie avec let me + infinitif sans to
\item À la fin de cette leçon, je sais formuler une demande polie avec could you et can you + base verbale
\item À la fin de cette leçon, je sais placer correctement please dans une demande polie
\item À la fin de cette leçon, je sais distinguer les niveaux de politesse (can you < could you < could you ... please)
\item À la fin de cette leçon, je sais éviter les erreurs fréquentes des francophones (to après let, traduction mot à mot de 'pouvez-vous')
\item À la fin de cette leçon, je sais répondre correctement à une demande polie (acceptation et refus poli)
\item À la fin de cette leçon, je sais utiliser ces structures dans un dialogue professionnel court (marché, bureau, banque)
\end{itemize}
\end{objectifs}

\begin{prerequis}
\begin{itemize}
\item Les verbes modaux can et could (Seconde)
\item L'infinitif avec et sans to (bare infinitive)
\item La structure de la phrase interrogative en anglais (inversion sujet-verbe)
\item Le vocabulaire de base de la vie économique (shop, market, bank, office)
\item Les formules de salutation et de politesse simples (hello, thank you, excuse me)
\end{itemize}
\end{prerequis}

\newpage

\coursec{Observation et découverte}

\courssub{Situation d'accroche : politeness first!}

Imagine you are at the Douala International Airport. A British tourist is struggling with a heavy suitcase. You want to help, but you hesitate: should you say \eng{Give me your bag!}? That sounds rude in English! In English-speaking countries like the United Kingdom, the United States or the South West Region of Cameroon, politeness is essential: people say \eng{Let me help you with your bag} or \eng{Could you open the window, please?}. Today, you will learn how to make polite requests and suggestions like a real English speaker.

En français, nous disons « Laisse-moi t'aider », « Pourrais-tu me montrer... ? », « Peux-tu baisser le prix, s'il te plaît ? ». L'anglais possède des outils équivalents, mais avec des mécanismes de formation différents : pas de subjonctif, pas de conditionnel compliqué, mais des structures très codifiées qu'il faut mémoriser telles quelles. Cette leçon va te les faire découvrir par l'observation, comme le fait un linguiste : d'abord un texte authentique, ensuite des hypothèses, et enfin la règle.

\textbf{Problématique de la leçon :} comment formuler une demande ou une suggestion de manière polie en anglais, sans jamais avoir l'air de donner un ordre ?

\courssub{Lecture du dialogue : At the market in Bamenda}

Lis le dialogue suivant à voix haute, en faisant attention à l'intonation montante des questions polies. La scène se déroule au marché principal de Bamenda, dans la région du Nord-Ouest : une cliente veut acheter du tissu (fabric) pour une cérémonie.

\begin{textebox}[At the market — Bamenda Main Market]
Customer: Excuse me, could you show me that fabric, please? The blue one on the top shelf.

Vendor: Of course, madam! Let me help you. Here it is. It is a beautiful wax print from Nigeria.

Customer: It is lovely. How much is it?

Vendor: It is 7,000 FCFA per piece.

Customer: Oh, that is a bit expensive. Can you reduce the price, please?

Vendor: Let me see... You are my first customer today. OK, 5,000 FCFA for you.

Customer: Thank you very much! Could you wrap it for me, please?

Vendor: Sure! Let me find a nice bag for you.

Customer: You are very kind. Have a good day!

Vendor: Thank you, madam. Come again!
\end{textebox}

\begin{definition}[Glossaire du dialogue]
\begin{itemize}
\item \cle{fabric} (nom) : le tissu, l'étoffe.
\item \cle{shelf} (nom, pluriel \eng{shelves}) : l'étagère. « on the top shelf » = sur l'étagère du haut.
\item \cle{wax print} (nom) : le pagne wax, tissu imprimé.
\item \cle{to reduce} (verbe) : réduire, baisser (le prix). Prononciation : « ri-diouce ».
\item \cle{to wrap} (verbe) : emballer, envelopper. Prononciation : « rap ».
\item \cle{kind} (adjectif) : gentil, aimable.
\end{itemize}
\end{definition}

\courssub{Repérage : trouve les phrases polies}

Relis le dialogue et souligne toutes les phrases qui expriment soit une \cle{demande} (le locuteur veut que l'autre fasse quelque chose), soit une \cle{suggestion} ou une \cle{offre d'aide} (le locuteur propose de faire quelque chose lui-même). Tu devrais trouver cinq énoncés :

\begin{exemplebox}[Les cinq énoncés polis du dialogue]
\begin{enumerate}
\item \eng{Could you show me that fabric, please?} « Pourrais-tu me montrer ce tissu, s'il te plaît ? »
\item \eng{Let me help you.} « Laisse-moi t'aider. »
\item \eng{Can you reduce the price, please?} « Peux-tu baisser le prix, s'il te plaît ? »
\item \eng{Let me see...} « Laisse-moi voir... »
\item \eng{Could you wrap it for me, please?} « Pourrais-tu l'emballer pour moi, s'il te plaît ? »
\end{enumerate}
\end{exemplebox}

\begin{methode}[Comment observer un énoncé de grammaire]
Pour analyser une phrase polie en anglais, procède toujours en trois étapes :
\begin{enumerate}
\item \textbf{Souligne le verbe principal} de la phrase (par exemple \eng{show}, \eng{help}, \eng{reduce}, \eng{wrap}).
\item \textbf{Identifie la structure qui précède} ce verbe : est-ce \eng{Let me}, \eng{Could you} ou \eng{Can you} ?
\item \textbf{Vérifie si le verbe principal porte « to » ou non} : dans toutes nos phrases, le verbe est à la \cle{base verbale} (infinitif sans \eng{to}). On dit \eng{Let me help}, jamais \eng{Let me to help}.
\end{enumerate}
\end{methode}

\courssub{Question guidée : qu'est-ce qui rend ces phrases polies ?}

Compare les deux colonnes du tableau suivant. À gauche, des phrases directes qui sonnent comme des ordres ; à droite, les versions polies du dialogue.

\begin{center}
\begin{tabularx}{\linewidth}{|X|X|X|}
\hline
\rowcolor{popblueL} Phrase directe (rude) & Version polie & Élément de politesse \\
\hline
Show me that fabric! & Could you show me that fabric, please? & \eng{could you} + \eng{please} \\
\hline
Help me! & Let me help you. & \eng{let me} (offre d'aide) \\
\hline
Reduce the price! & Can you reduce the price, please? & \eng{can you} + \eng{please} \\
\hline
Wrap it for me! & Could you wrap it for me, please? & \eng{could you} + \eng{please} \\
\hline
\end{tabularx}
\end{center}

\textbf{Questions de réflexion (à discuter en classe) :}
\begin{itemize}
\item Quels mots reviennent au début des demandes polies ? (\eng{Could you...?}, \eng{Can you...?})
\item Quel petit mot magique apparaît à la fin des demandes ? (\eng{please})
\item Dans les suggestions, quel mot précède le pronom \eng{me} ? (\eng{Let})
\item Le verbe principal porte-t-il \eng{to}, \eng{-s} ou \eng{-ing} ? (Non : il reste à la base verbale.)
\end{itemize}

\begin{attention}[Demande polie ou ordre ?]
Ne confonds jamais une demande polie avec un ordre. \eng{Open the door!} est un \cle{ordre} (impératif) : il peut paraître brutal ou impoli. \eng{Could you open the door, please?} est une \cle{demande polie} : c'est la forme attendue dans la vie quotidienne anglophone, même en famille ou entre amis. En anglais, l'impératif direct est réservé aux urgences (\eng{Run!}), aux panneaux (\eng{Keep off the grass}) ou aux consignes d'examen.
\end{attention}

\courssub{Classement : suggestions ou demandes ?}

Classons maintenant les cinq énoncés relevés en deux catégories grammaticales. Observe bien le sujet : quand le locuteur propose de faire l'action \emph{lui-même}, il utilise \eng{Let me...} ; quand il demande à \emph{l'autre personne} de faire l'action, il utilise \eng{Could you...?} ou \eng{Can you...?}.

\begin{center}
\begin{tabularx}{\linewidth}{|X|X|}
\hline
\rowcolor{popgreenL} Suggestions / offres d'aide : Let me + V & Demandes polies : Could/Can you + V \\
\hline
Let me help you. & Could you show me that fabric, please? \\
\hline
Let me see... & Can you reduce the price, please? \\
\hline
Let me find a nice bag for you. & Could you wrap it for me, please? \\
\hline
\end{tabularx}
\end{center}

\begin{popfigure}
\begin{center}\tcbox[colback=popink,colframe=popink,coltext=white,arc=9pt,boxsep=4pt,left=6pt,right=6pt]{\bfseries\small Polite English}\end{center}
\vspace{-2pt}
\begin{tcbraster}[raster columns=2,raster equal height=rows,raster column skip=6pt,raster row skip=6pt,enhanced,blankest]
\begin{tcolorbox}[enhanced,colback=white,colframe=popgreen!70,boxrule=0.9pt,arc=3pt,title={Suggestions},fonttitle=\bfseries\scriptsize,coltitle=white,colbacktitle=popgreen!70,left=4pt,right=4pt,top=2pt,bottom=2pt,boxsep=1pt,toptitle=1.5pt,bottomtitle=1.5pt,halign=left]
\begin{itemize}[nosep,leftmargin=*,itemsep=1pt,font=\scriptsize]
\item Let me + V \emph{Let me help you.}
\end{itemize}
\end{tcolorbox}
\begin{tcolorbox}[enhanced,colback=white,colframe=poporange!75,boxrule=0.9pt,arc=3pt,title={Requests},fonttitle=\bfseries\scriptsize,coltitle=white,colbacktitle=poporange!75,left=4pt,right=4pt,top=2pt,bottom=2pt,boxsep=1pt,toptitle=1.5pt,bottomtitle=1.5pt,halign=left]
\begin{itemize}[nosep,leftmargin=*,itemsep=1pt,font=\scriptsize]
\item Could/Can you + V ... please? \emph{Could you wait,} \emph{please?}
\end{itemize}
\end{tcolorbox}
\end{tcbraster}

\legende{Carte mentale : les deux grandes familles de la politesse anglaise — la suggestion avec \eng{let me} et la demande avec \eng{could you / can you}.}
\end{popfigure}

\courssub{Hypothèses des élèves : à toi de formuler la règle !}

Avant la leçon magistrale, essaie de compléter ces hypothèses avec ton voisin. C'est la démarche d'observation guidée : le cerveau retient mieux une règle qu'il a d'abord devinée.

\begin{enumerate}
\item Après \eng{Let me}, le verbe se met à \dots (l'infinitif avec \eng{to} ? la base verbale sans \eng{to} ? la forme en \eng{-ing} ?)
\item \eng{Could} est le passé de \eng{can}. Dans une demande polie, lequel des deux est le \emph{plus} poli ? \dots
\item Le mot \eng{please} se place plutôt au \dots de la phrase.
\item Pour traduire « Laisse-moi voir », l'anglais utilise le verbe \eng{to let} suivi de \dots puis de \dots
\end{enumerate}

\begin{exoresolu}[Premier pas : transforme l'ordre en demande polie]
Énoncé : transforme cet ordre brutal en demande polie adressée à un vendeur du marché : \eng{Give me a discount!}
\tcblower
Solution détaillée :
\begin{enumerate}
\item On souligne le verbe : \eng{give}. Il restera à la base verbale, sans \eng{to}.
\item On choisit l'outil de politesse pour une demande : \eng{Could you} (plus poli) ou \eng{Can you} (courant).
\item On place \eng{please} à la fin de la phrase.
\end{enumerate}
Réponse attendue : \eng{Could you give me a discount, please?} « Pourriez-vous me faire une réduction, s'il vous plaît ? »
Variante acceptée : \eng{Can you give me a discount, please?}
\end{exoresolu}

\begin{savaistu}[Please et thank you : les mots magiques des Britanniques]
Au Royaume-Uni, dire \eng{please} et \eng{thank you} est presque obligatoire dans chaque échange : au magasin, au restaurant, dans le bus, et même entre amis proches ou en famille. Oublier \eng{please} est perçu comme impoli, même si ta phrase est grammaticalement correcte ! Les Britanniques appellent ces mots les « magic words ». Un enfant britannique qui demande \eng{I want water!} entendra souvent sa mère répondre : \eng{What's the magic word?} — et l'enfant doit ajouter \eng{please}. Au Cameroun anglophone, dans les villes comme Buea ou Bamenda, la même culture de la politesse verbale existe : un client qui dit simplement \eng{Give me bread!} au marché passera pour mal élevé.
\end{savaistu}

\begin{exercice}[À toi de jouer : observe et classe]
Voici quatre nouvelles phrases entendues au marché de Mokolo à Yaoundé. Pour chacune : a) souligne le verbe principal ; b) dis s'il s'agit d'une suggestion (S) ou d'une demande polie (D) ; c) traduis-la en français.
\begin{enumerate}
\item \eng{Let me carry your basket, madam.}
\item \eng{Can you weigh these tomatoes, please?}
\item \eng{Could you speak slowly, please?}
\item \eng{Let me check the price for you.}
\end{enumerate}
\end{exercice}

\begin{corrige}[Exercice : observe et classe]
\begin{enumerate}
\item Verbe : \eng{carry}. Suggestion (S) : le locuteur propose de porter lui-même. « Laissez-moi porter votre panier, madame. »
\item Verbe : \eng{weigh} (peser). Demande polie (D). « Pouvez-vous peser ces tomates, s'il vous plaît ? »
\item Verbe : \eng{speak}. Demande polie (D). « Pourriez-vous parler lentement, s'il vous plaît ? »
\item Verbe : \eng{check} (vérifier). Suggestion (S). « Laissez-moi vérifier le prix pour vous. »
\end{enumerate}
Remarque : dans les quatre cas, le verbe est à la base verbale, sans \eng{to} ni \eng{-s}. C'est exactement la règle que nous allons énoncer dans la section suivante.
\end{corrige}

\begin{aretenir}[Ce que l'observation nous a déjà appris]
\begin{itemize}
\item Une \cle{suggestion} ou une offre d'aide se construit avec \eng{Let me} + base verbale : \eng{Let me help you.}
\item Une \cle{demande polie} se construit avec \eng{Could you} ou \eng{Can you} + base verbale, souvent suivie de \eng{please} : \eng{Could you open the door, please?}
\item \eng{Could you} est plus poli et plus formel que \eng{Can you}.
\item Le verbe qui suit ces structures ne prend jamais \eng{to}, jamais \eng{-s}, jamais \eng{-ing}.
\end{itemize}
\end{aretenir}

\coursec{Lesson 14 — Grammar: Polite requests and suggestions}

\courssub{Observation et découverte : un dialogue au marché de Mokolo}

Pour entrer dans la leçon, écoutons une interaction typique dans un contexte économique camerounais. Ce dialogue met en scène un client et un vendeur de téléphones à Douala. Repérez les formules utilisées pour offrir de l'aide, proposer un service ou suggérer une action commune.

\begin{textebox}[At the phone shop — Douala]
\textbf{Customer:} Excuse me, how much is this phone?

\textbf{Seller:} Let me check the price for you. Ah, it is 45,000 francs.

\textbf{Customer:} That is expensive. Let me think about it.

\textbf{Seller:} Let me show you a cheaper model. This one is 30,000 francs.

\textbf{Customer:} Good. Let me see if I have enough money... Yes, I do.

\textbf{Seller:} Perfect. Let me wrap it for you. And let's not forget your receipt!

\textbf{Customer:} Thank you very much.

\textbf{Seller:} You are welcome. Come again!
\end{textebox}

\textbf{Glossaire :}

\begin{center}
\begin{tabularx}{\linewidth}{|l|l|X|}
\hline
\rowcolor{popblueL} English & Nature & Français \\
\hline
to check & verbe & vérifier \\
\hline
to wrap & verbe & emballer \\
\hline
receipt & nom & reçu \\
\hline
cheap / cheaper & adjectif & bon marché / moins cher \\
\hline
to think about it & expression & y réfléchir \\
\hline
expensive & adjectif & cher \\
\hline
enough & déterminant/adverbe & assez \\
\hline
model & nom & modèle \\
\hline
\end{tabularx}
\end{center}

\courssub{Analyse guidée}

Observez les phrases mises en évidence dans le dialogue : \eng{Let me check}, \eng{Let me think}, \eng{Let me show}, \eng{Let me see}, \eng{Let me wrap}, \eng{let's not forget}.

\begin{enumerate}
\item Quel mot suit systématiquement \eng{let me} ? Un verbe. Sous quelle forme ? \textbf{La forme de base (infinitif sans \eng{to})}.
\item Quelle est la différence entre \eng{Let me help you} et \eng{Let's go to the bank} ? \eng{Let me} engage \textbf{seulement le locuteur} (offre d'aide) ; \eng{Let's} (\eng{let us}) engage \textbf{le locuteur ET l'interlocuteur} (suggestion commune).
\item Traduisez mentalement : \eng{Let me check the price} $\rightarrow$ « Laisse-moi vérifier le prix » ; \eng{Let's not forget} $\rightarrow$ « N'oublions pas ».
\end{enumerate}

\coursec{Règle 1 : Let me + base verbale (offres d'aide, suggestions personnelles)}

\courssub{Structure et formation}

\begin{propriete}[La structure « let me + base verbale »]
Après \cle{let me}, le verbe est toujours à l'\textbf{infinitif SANS \eng{to}} (appelé \emph{bare infinitive} ou \textbf{base verbale}) :

\begin{center}
\textbf{Let me + VERBE (base form) + complément}
\end{center}

\eng{Let me help you.} = Laisse-moi t'aider. \hfill (Jamais \eng{Let me to help you})

\eng{Let me open the door for you.} = Laisse-moi t'ouvrir la porte.
\end{propriete}

\begin{popfigure}
\begin{tikzpicture}[
box/.style={draw=popblue, fill=popblueL, rounded corners=3pt, align=center, minimum height=1.1cm, font=\small, text width=2.2cm},
verbe/.style={draw=poporange, fill=poporangeL, rounded corners=3pt, align=center, minimum height=1.1cm, font=\small, text width=2.2cm},
comp/.style={draw=popgreen, fill=popgreenL, rounded corners=3pt, align=center, minimum height=1.1cm, font=\small, text width=3cm},
plus/.style={font=\large\bfseries, text=popdark}
]
\node[box, align=center] (a) at (0,0){\textbf{Let me}\\ \emph{(invariable)}};
\node[plus] at (2.5,0) {+};
\node[verbe, align=center] (b) at (4.5,0){\textbf{carry}\\ \emph{verbe à la base}};
\node[plus] at (7.0,0) {+};
\node[comp, align=center] (c) at (9.5,0){\textbf{your bag for you}\\ \emph{complément}};
\node[below=0.7cm of b, align=center, text=popdark, font=\small] {Let me carry your bag for you.\\ « Laisse-moi porter ton sac pour toi. »};
\end{tikzpicture}
\legende{Schéma de structure : let me + base verbale + complément}
\end{popfigure}

\begin{attention}[Point de vigilance : « let me » est invariable]
\cle{Let me} ne se conjugue \textbf{jamais}. Il ne prend pas de \eng{-s} à la 3\textsuperscript{e} personne, car il exprime toujours ce que \emph{je} (le locuteur) propose de faire. On ne dit pas \eng{He lets me help you} pour exprimer une offre de sa part. Pour une tierce personne, on change de structure : \eng{He wants to help you.} / \eng{He offers to help you.}
\end{attention}

\courssub{Les trois emplois principaux}

\textbf{1. Offrir son aide spontanément} (réaction à une difficulté visible).

\begin{exemplebox}[Offrir son aide]
\eng{Let me help you with your luggage.} = Laisse-moi t'aider avec tes bagages.

\eng{Let me carry that heavy box, Sir.} = Laissez-moi porter ce carton lourd, Monsieur.

\eng{You look lost. Let me show you the way.} = Tu as l'air perdu. Laisse-moi te montrer le chemin.
\end{exemplebox}

\textbf{2. Proposer de prendre en charge une tâche} (contexte service, achat, bureau).

\begin{exemplebox}[Proposer un service]
\eng{Let me pay for the taxi.} = Laisse-moi payer le taxi.

\eng{Let me explain the procedure.} = Laisse-moi t'expliquer la procédure.

\eng{Let me get you a glass of water.} = Laisse-moi t'apporter un verre d'eau.

\eng{Let me book the ticket for you.} = Laisse-moi réserver le billet pour toi.
\end{exemplebox}

\textbf{3. Annoncer une initiative polie} (commerce, banque, administration : « je vais faire ceci pour vous »).

\begin{exemplebox}[Initiative polie en contexte économique]
\eng{Let me check the availability.} = Laisse-moi vérifier la disponibilité.

\eng{Let me ask my manager.} = Laisse-moi demander à mon responsable.

\eng{Let me give you a 10\% discount.} = Laisse-moi vous accorder 10\% de remise.

\eng{Let me write down your contact details.} = Laisse-moi noter vos coordonnées.
\end{exemplebox}

\courssub{Extension : « let's » pour suggérer une action commune}

Quand la suggestion inclut \textbf{le locuteur ET son/ses interlocuteur(s)}, on utilise \cle{let's} (contraction de \emph{let us}) + base verbale. En français : impératif 1\textsuperscript{re} personne du pluriel (« allons », « faisons », « ne ... pas »).

\begin{propriete}[La structure « let's + base verbale »]
\begin{center}
\textbf{Let's + VERBE (base form) + complément} = suggestion d'action commune
\end{center}

Forme négative : \textbf{Let's not + VERBE (base form)} = suggestion de ne pas faire ensemble.
\end{propriete}

\begin{exemplebox}[Suggestions avec « let's »]
\eng{Let's go to the bank now.} = Allons à la banque maintenant.

\eng{Let's compare prices before buying.} = Comparons les prix avant d'acheter.

\eng{Let's not waste time.} = Ne perdons pas de temps.

\eng{Let's not argue about the price.} = Ne nous disputons pas pour le prix.

\eng{Let's meet at the entrance of the market.} = Retrouvons-nous à l'entrée du marché.
\end{exemplebox}

\begin{attention}[La négation avec « let's » : piège classique]
La négation se place \textbf{AVANT} le verbe : \eng{Let's \textbf{not} waste time.}
\begin{itemize}
\item Interdit : \eng{Let's don't waste time.} (\eng{don't} n'existe pas après \eng{let's}).
\item Interdit : \eng{Let's to not waste time.} (pas de \eng{to}).
\end{itemize}
Le verbe reste à la base après \eng{not}.
\end{attention}

\courssub{Comment répondre à une offre (\eng{let me}) ou une suggestion (\eng{let's})}

\begin{center}
\begin{tabularx}{\linewidth}{|X|X|}
\hline
\rowcolor{popblueL} Réponse en anglais & Contexte / Traduction \\
\hline
Yes, please. / Yes, thank you. & Accepter une offre (\eng{Let me help}). \\
\hline
That's very kind of you. & Accepter avec gratitude (formel). \\
\hline
Good idea! / Great idea! & Accepter une suggestion (\eng{Let's go}). \\
\hline
No, thank you. I can manage. & Refuser poliment une offre. \\
\hline
I'd rather not, if you don't mind. & Refuser une suggestion (\eng{Let's...}) poliment. \\
\hline
\end{tabularx}
\end{center}

\coursec{Règle 2 : Could you / Can you + base verbale (demandes polies)}

\courssub{Observation : une scène à la banque}

Lisons cet échange à la banque (BICEC, Yaoundé). Le client a besoin de l'aide du guichetier.

\begin{textebox}[At the bank counter — Yaoundé]
\textbf{Customer:} Good morning. \textbf{Could you help me fill in this form, please?}

\textbf{Clerk:} Good morning. \textbf{Certainly.} Let me get a pen for you. \textbf{Can you sign here, please?}

\textbf{Customer:} \textbf{Of course.} \textbf{Could you also check my balance?}

\textbf{Clerk:} \textbf{Yes, I can.} One moment... Your balance is 125,000 francs.

\textbf{Customer:} Thank you very much. \textbf{Would you mind stamping the paper?}

\textbf{Clerk:} \textbf{Not at all.} Here you are.
\end{textebox}

\textbf{Glossaire :}

\begin{center}
\begin{tabularx}{\linewidth}{|l|l|X|}
\hline
\rowcolor{popblueL} English & Nature & Français \\
\hline
to fill in / out & verbe phrasal & remplir (un formulaire) \\
\hline
balance & nom & solde (bancaire) \\
\hline
to stamp & verbe & tamponner, estampiller \\
\hline
certainly / of course / not at all & adverbes / locutions & bien sûr / certainement / il n'y a pas de quoi \\
\hline
would you mind + -ing & structure & est-ce que cela vous dérange de... \\
\hline
\end{tabularx}
\end{center}

\courssub{Analyse guidée}

\begin{enumerate}
\item Repérez les demandes du client : \eng{Could you help me...?}, \eng{Could you also check...?}, \eng{Would you mind stamping...?}.
\item Repérez la demande du guichetier : \eng{Can you sign here...?}.
\item Quelle forme verbale suit \eng{Could you}, \eng{Can you} ? \textbf{La base verbale (infinitif sans \eng{to})}.
\item Quelle est la nuance entre \eng{Can you} et \eng{Could you} ? \eng{Could you} est \textbf{plus poli, plus formel} (conditionnel de politesse). \eng{Can you} est standard, neutre, souvent utilisé entre collègues ou par un supérieur/employé vers un client.
\end{enumerate}

\courssub{La règle : formes et emplois}

\begin{propriete}[Demandes polies : « Could you / Can you + base verbale »]
\begin{center}
\textbf{Can / Could + YOU + VERBE (base form) + complément ?}
\end{center}

\begin{itemize}
\item \textbf{Can you... ?} : Demande standard, neutre. « Peux-tu / Pouvez-vous... ? »
\item \textbf{Could you... ?} : Demande plus polie, plus formelle. « Pourriez-vous... ? » (recommandé avec un inconnu, un supérieur, dans l'administration).
\item \textbf{Would you mind + VERBE-\eng{ing}... ?} : Très poli, « Est-ce que cela vous dérange de... ? » (ex. \eng{Would you mind opening the window?}).
\end{itemize}
\end{propriete}

\begin{popfigure}
\begin{tikzpicture}[
modal/.style={draw=poppurple, fill=poppurpleL, rounded corners=3pt, align=center, minimum height=1cm, font=\small, text width=1.8cm},
subj/.style={draw=popblue, fill=popblueL, rounded corners=3pt, align=center, minimum height=1cm, font=\small, text width=1.2cm},
verb/.style={draw=poporange, fill=poporangeL, rounded corners=3pt, align=center, minimum height=1cm, font=\small, text width=2cm},
comp/.style={draw=popgreen, fill=popgreenL, rounded corners=3pt, align=center, minimum height=1cm, font=\small, text width=3.5cm},
qmark/.style={font=\large\bfseries, text=popdark},
plus/.style={font=\large\bfseries, text=popdark}
]
\node[modal] (m) at (0,0) {\textbf{Could / Can}};
\node[plus] at (2.1,0) {+};
\node[subj] (s) at (3.8,0) {\textbf{you}};
\node[plus] at (5.2,0) {+};
\node[verb, align=center] (v) at (7.0,0){\textbf{help}\\ \emph{base verbale}};
\node[plus] at (9.3,0) {+};
\node[comp, align=center] (c) at (11.5,0){\textbf{me with this form?}\\ \emph{complément}};
\node[qmark] at (14.5,0) {?};
\node[below=0.7cm of v, align=center, text=popdark, font=\small] {Could you help me with this form?\\ « Pourriez-vous m'aider avec ce formulaire ? »};
\end{tikzpicture}
\legende{Schéma de structure : Could/Can you + base verbale}
\end{popfigure}

\courssub{Le mot « please » : position et intonation}

\cle{Please} adoucit la demande. Trois positions possibles :
\begin{enumerate}
\item \textbf{En fin de phrase} (le plus courant, neutre) : \eng{Could you help me, please?}
\item \textbf{Au début} (insistance, formalité) : \eng{Please, could you wait here?}
\item \textbf{Entre le modal et le sujet} (très formel, écrit administratif) : \eng{Could you please send the letter?}
\end{enumerate}

\begin{attention}[Piège : « Can you » vs « Could you » à l'écrit et à l'oral]
En classe et aux examens, \textbf{privilégiez \eng{Could you} pour toute demande adressée à un adulte, un inconnu, un professeur, un guichetier}. \eng{Can you} est acceptable entre amis ou pour une demande simple à un commerçant connu. \eng{Will you} / \eng{Would you} existent aussi (\eng{Would you} = poli, \eng{Will you} = insistant/ordonne).
\end{attention}

\courssub{Répondre à une demande (\eng{Can/Could you})}

\begin{center}
\begin{tabularx}{\linewidth}{|X|X|}
\hline
\rowcolor{popblueL} Réponse positive & Réponse négative (polie) \\
\hline
Certainly. / Of course. / Sure. & I'm afraid I can't. / I'm sorry, I can't. \\
\hline
Yes, I can. / Yes, of course. & I'd like to, but... (reason). \\
\hline
No problem. / With pleasure. & Not right now, maybe later. \\
\hline
\end{tabularx}
\end{center}

\begin{savaistu}[« Would you mind » + \eng{-ing} : la politesse maximale]
\eng{Would you mind opening the window?} = « Est-ce que cela vous dérange d'ouvrir la fenêtre ? »
\begin{itemize}
\item Réponse positive (j'accepte de le faire) : \eng{Not at all.} / \eng{Of course not.} (Attention : en français on dit « Bien sûr », en anglais on nie le dérangement : « Pas du tout »).
\item Réponse négative (je refuse) : \eng{I'm sorry, but I'd rather not.} / \eng{Well, actually I do mind.}
\end{itemize}
C'est une structure de niveau B1/B2, très prisée au Baccalauréat pour l'expression écrite formelle (lettres de motivation, réclamations).
\end{savaistu}

\coursec{Comparaison français / anglais et pièges des francophones}

\courssub{Tableau récapitulatif des correspondances}

\begin{center}
\begin{tabularx}{\linewidth}{|X|X|X|}
\hline
\rowcolor{popblueL} Français & English & Remarque cruciale \\
\hline
Laisse-moi t'aider. & Let me help you. & \textbf{Pas de \eng{to}} après \eng{let me}. \\
\hline
Laisse-moi payer le taxi. & Let me pay \textbf{for} the taxi. & \eng{Pay} qqn \textbf{mais} \eng{pay for} qqch. \\
\hline
Allons à la banque. & Let's go to the bank. & \eng{Let's} = \eng{let us}. \\
\hline
Ne gaspillons pas d'argent. & Let's not waste money. & \eng{Not} avant le verbe, pas \eng{don't}. \\
\hline
Peux-tu m'aider ? (familier/standard) & Can you help me? & \eng{Can} = capacité/permission simple. \\
\hline
Pourriez-vous m'aider ? (poli/formel) & \textbf{Could you help me?} & \textbf{Forme à privilégier en contexte pro/admin.} \\
\hline
Est-ce que cela vous dérange d'aider ? & Would you mind helping me? & \textbf{Verbe en \eng{-ing}} après \eng{mind}. \\
\hline
\end{tabularx}
\end{center}

\courssub{Les 5 pièges majeurs des francophones}

\begin{attention}[Piège n°1 : Le « to » parasite après \eng{let me} / \eng{let's}]
\eng{FAUX : Let me to help you.} \quad \eng{FAUX : Let's to go.}
\eng{BON : Let me help you.} \quad \eng{BON : Let's go.}
\emph{Astuce} : Le « t' » de « laisse-moi t'aider » est un pronom objet (\emph{te}), pas la préposition \eng{to}.
\end{attention}

\begin{attention}[Piège n°2 : Le « don't » parasite avec \eng{let's not}]
\eng{FAUX : Let's don't go.} \quad \eng{FAUX : Let's not to go.}
\eng{BON : Let's not go.}
\end{attention}

\begin{attention}[Piège n°3 : Le « -s » de la 3\textsuperscript{e} personne sur le verbe après \eng{let me}]
\eng{FAUX : Let me helps you.} \quad \eng{FAUX : Let me carries the bag.}
\eng{BON : Let me help you.} \quad \eng{BON : Let me carry the bag.}
\emph{Rappel} : \eng{let me} est invariable ; le verbe qui suit est une base verbale (infinitif sans \eng{to}), jamais conjugué.
\end{attention}

\begin{attention}[Piège n°4 : Confondre \eng{Can you} (demande) et \eng{You can} (permission/affirmation)]
\eng{Can you open the door?} = Peux-tu / Pourriez-vous ouvrir la porte ? (Question/Requête).
\eng{You can open the door.} = Tu peux / Vous pouvez ouvrir la porte. (Affirmation/Permission donnée).
\emph{À l'oral}, l'intonation montante distingue la question ; \emph{à l'écrit}, le point d'interrogation et l'inversion (modal + sujet) sont obligatoires.
\end{attention}

\begin{attention}[Piège n°5 : Oublier le \eng{-ing} après \eng{Would you mind}]
\eng{FAUX : Would you mind to help me?} \quad \eng{FAUX : Would you mind help me?}
\eng{BON : Would you mind \textbf{helping} me?}
\emph{Règle} : \eng{mind} (comme \eng{enjoy, avoid, finish, suggest}) est suivi du gérondif (\eng{V-ing}).
\end{attention}

\coursec{Entraînement guidé et usage en contexte économique}

\courssub{Exercice résolu 1 : Transformation (Let me / Let's)}

\begin{exoresolu}[Transformer en utilisant \eng{let me} ou \eng{let's}]
Énoncé — Reformulez les phrases suivantes en utilisant la structure indiquée.

1. I want to help you with your bags. $\rightarrow$ \eng{Let me...}
2. We should go to the market now. $\rightarrow$ \eng{Let's...}
3. I propose to pay the bill. $\rightarrow$ \eng{Let me...}
4. We shouldn't spend all our money. $\rightarrow$ \eng{Let's not...}
5. I will check the timetable for you. $\rightarrow$ \eng{Let me...}

\tcblower
Solution détaillée —

1. \textbf{Let me help you with your bags.} (Offre d'aide : \eng{let me} + base).
2. \textbf{Let's go to the market now.} (Suggestion commune : \eng{let's} + base).
3. \textbf{Let me pay the bill.} (Proposition de service : \eng{let me} + base ; attention : \eng{pay the bill}, pas \eng{pay for the bill} ici car « payer l'addition » = \eng{pay the bill}).
4. \textbf{Let's not spend all our money.} (Négation suggestion : \eng{let's not} + base).
5. \textbf{Let me check the timetable for you.} (Initiative polie : \eng{let me} + base).
\end{exoresolu}

\courssub{Exercice résolu 2 : Choix du modal (Can / Could / Would you mind)}

\begin{exoresolu}[Compléter avec la forme appropriée]
Énoncé — Choisissez entre \eng{Can you}, \eng{Could you}, \eng{Would you mind} + \eng{-ing}. Adaptez le verbe.

1. \ldots\ \ldots\ (pass) me the salt, please? (Contexte : dîner en famille, standard).
2. \ldots\ \ldots\ (open) the window? (Contexte : bureau, à un collègue qu'on connaît peu, poli).
3. \ldots\ \ldots\ (wait) a moment? (Contexte : guichet de poste, très formel).
4. \ldots\ \ldots\ (help) me carry this table? (Contexte : ami proche, demande directe).

\tcblower
Solution détaillée —

1. \textbf{Can you pass} me the salt, please? (Standard, familier/amical).
2. \textbf{Could you open} the window? (Poli, recommandé pour connaissance superficielle).
3. \textbf{Would you mind waiting} a moment? (Très formel, administration ; \eng{waiting} = gérondif).
4. \textbf{Can you help} me carry this table? (Direct, entre amis). \eng{Could you help} est aussi possible (plus poli).
\end{exoresolu}

\courssub{Exercice résolu 3 : Traduction thématique (Contexte économique)}

\begin{exoresolu}[Traduire en anglais]
Énoncé —

1. Laisse-moi vérifier le taux de change pour vous. (Banque)
2. Pourriez-vous signer ce document, s'il vous plaît ? (Administration)
3. Allons comparer les prix au marché Central. (Achats)
4. Ne perdons pas de temps, le guichet ferme à 15h. (Urgence)
5. Est-ce que cela vous dérange de m'attendre cinq minutes ? (Politesse maximale)

\tcblower
Solution détaillée —

1. \textbf{Let me check the exchange rate for you.}
2. \textbf{Could you sign this document, please?} (Ou \eng{Would you mind signing this document?})
3. \textbf{Let's compare prices at the Central Market.}
4. \textbf{Let's not waste time, the counter closes at 3 p.m.}
5. \textbf{Would you mind waiting for me five minutes?}
\end{exoresolu}

\courssub{Mise en situation orale : Jeu de rôle (Speaking)}

\begin{experience}[Role-play : Au marché / À la banque / Au cybercafé]
\textbf{Consigne :} Par binômes. Préparez un dialogue de 1 minute environ en utilisant \textbf{au moins 3 structures} étudiées (\eng{let me}, \eng{let's}, \eng{could you}, \eng{can you}, \eng{would you mind}).

\textbf{Scénario A (Marché Mokolo) :} Un client achète du tissu. Le vendeur propose de montrer la qualité, de couper, d'emballer. Le client suggère d'aller boire un jus après.
\textbf{Scénario B (Banque BICEC) :} Un client veut ouvrir un compte / faire un virement. Le guichetier demande de remplir un formulaire, de signer, de présenter une pièce d'identité. Le client demande une attestation.
\textbf{Scénario C (Cybercafé, Buea) :} Un étudiant veut imprimer son CV. Le gérant propose de vérifier le format, de le mettre sur clé USB. L'étudiant demande de scanner un document.

\textbf{Critères d'évaluation :} Correction des structures (base verbale, \eng{-ing} après \eng{mind}), intonation montante pour les questions, politesse (\eng{please}, \eng{thank you}), vocabulaire économique (receipt, balance, print, scan, discount, change).
\end{experience}

\courssub{Production écrite guidée (Writing)}

\begin{methode}[Rédiger un email formel de demande (niveau B1/B2)]
\textbf{Situation :} Vous êtes stagiaire dans une entreprise à Douala. Vous écrivez à votre tuteur (Mr. Ngando) pour lui demander de relire votre rapport de stage et de vous donner un créneau pour un entretien de fin de stage.

\textbf{Étapes :}
\begin{enumerate}
\item \textbf{Objet clair} : \eng{Subject: Request for feedback and meeting — Internship Report}.
\item \textbf{Formule d'appel} : \eng{Dear Mr. Ngando,}
\item \textbf{Contexte} : \eng{I have finished writing my internship report.}
\item \textbf{Demande 1 (polie, formelle)} : Utilisez \eng{Could you please...} ou \eng{Would you mind...} + \eng{-ing}. Ex. \eng{Could you please review it when you have time?}
\item \textbf{Demande 2 (proposition de rdv)} : Utilisez \eng{Could you let me know...} ou \eng{Let me know...}. Ex. \eng{Could you let me know a convenient time for a meeting?}
\item \textbf{Offre d'aide (optionnel)} : \eng{Let me know if you need a printed copy.}
\item \textbf{Formule de politesse finale} : \eng{Thank you for your time and guidance.} / \eng{Yours sincerely,} / \eng{Your name}.
\end{enumerate}

\end{methode}

\begin{exemplebox}[Modèle d'email]
\textbf{Subject:} Request for feedback — Internship Report

\textbf{Dear Mr. Ngando,}

I have completed my internship report. \textbf{Could you please review it} at your earliest convenience? \textbf{Would you mind letting me know} a suitable time for a debriefing meeting next week?

\textbf{Let me know} if you would like a printed copy.

Thank you for your guidance.

\textbf{Yours sincerely,}
\textbf{Jean Dupont}
\end{exemplebox}


\begin{exercice}[Entraînement complet — À rendre / À faire en classe]

\textbf{Partie A : Correction d'erreurs (Grammaire)}
Corrigez les phrases suivantes (une erreur par phrase).
1. Let me to help you with your homework.
2. Let's don't argue about the price.
3. Could you to open the door, please?
4. Would you mind to wait outside?
5. Let me carries this box for you, Madam.
6. Can you helping me fill this form?

\textbf{Partie B : Transformation (Let me / Let's / Could you)}
Réécrivez les phrases en utilisant la structure entre parenthèses.
1. I want to pay for the coffee. (\eng{Let me})
2. Shall we go to the supermarket? (\eng{Let's})
3. Do you want me to check the oil? (\eng{Let me})
4. I suggest we don't buy this car. (\eng{Let's not})
5. Is it possible for you to send the email? (\eng{Could you})

\textbf{Partie C : Traduction (Contexte camerounais)}
1. Laisse-moi te montrer où se trouve le guichet automatique.
2. Pourriez-vous me donner un reçu, s'il vous plaît ?
3. Allons négocier le prix du cacao chez l'acheteur.
4. Ne gaspillons pas notre salaire dans des futilités.
5. Est-ce que cela vous dérange de baisser la musique ? (Au cybercafé)

\textbf{Partie D : Expression libre}
Vous êtes vendeur dans une boutique de chaussures à Bamenda. Un client hésite entre deux paires. Écrivez un court dialogue (6-8 répliques) où vous utilisez \eng{let me} (offre), \eng{let's} (suggestion), \eng{could you} (demande au client), et une réponse polie du client.
\end{exercice}

\begin{corrige}[Exercice « Entraînement complet »]

\textbf{Partie A :}
1. \textbf{Let me help} you... (pas de \eng{to}).
2. \textbf{Let's not argue}... (pas de \eng{don't}).
3. \textbf{Could you open} the door... (pas de \eng{to} après modal).
4. \textbf{Would you mind waiting} outside? (\eng{mind} + \eng{-ing}).
5. \textbf{Let me carry} this box... (base verbale, pas de \eng{-s}).
6. \textbf{Can you help} me fill this form? (base verbale après \eng{can} ; \eng{help fill} ou \eng{help to fill}, mais pas \eng{helping}).

\textbf{Partie B :}
1. \textbf{Let me pay for the coffee.}
2. \textbf{Let's go to the supermarket.}
3. \textbf{Let me check the oil.}
4. \textbf{Let's not buy this car.}
5. \textbf{Could you send the email?}

\textbf{Partie C :}
1. \textbf{Let me show you where the ATM is.} (ou \eng{cash machine} / \eng{cashpoint}).
2. \textbf{Could you give me a receipt, please?}
3. \textbf{Let's negotiate the cocoa price with the buyer.}
4. \textbf{Let's not waste our salary on trivialities.} (ou \eng{frivolities}).
5. \textbf{Would you mind turning the music down?} (ou \eng{lowering the volume}).

\textbf{Partie D :} (Exemple de production attendue)
\textbf{Seller:} This leather pair is very durable. \textbf{Let me get your size.} \\
\textbf{Customer:} Thanks. They are a bit tight. \\
\textbf{Seller:} \textbf{Let me bring} the larger size. \textbf{Let's compare} them side by side. \\
\textbf{Customer:} Good idea. The brown ones look better. \\
\textbf{Seller:} \textbf{Could you try} walking with them? \\
\textbf{Customer:} \textbf{Certainly.} They are comfortable. I'll take the brown ones. \\
\textbf{Seller:} \textbf{Let me wrap} them for you.
\end{corrige}

\begin{aretenir}[Synthèse de la leçon 14 : Demandes et suggestions polies]
\begin{itemize}
\item \textbf{Let me + base verbale} : \textbf{Offre d'aide / Service / Initiative} (le locuteur agit). \eng{Let me help you.} \checkmark \quad \eng{Let me to help} \cross
\item \textbf{Let's + base verbale} : \textbf{Suggestion commune} (locuteur + interlocuteur). \eng{Let's go.} \checkmark
\item \textbf{Let's not + base verbale} : \textbf{Négation suggestion}. \eng{Let's not wait.} \checkmark \quad \eng{Let's don't wait} \cross
\item \textbf{Can you + base verbale} : \textbf{Demande standard/neutre}. \eng{Can you pass the salt?}
\item \textbf{Could you + base verbale} : \textbf{Demande polie/formelle} (examens, admin, inconnus). \eng{Could you sign here?} \checkmark
\item \textbf{Would you mind + V-\eng{ing}} : \textbf{Demande très polie}. \eng{Would you mind waiting?} \checkmark \quad \eng{Would you mind to wait?} \cross
\item \textbf{Please} : En fin de phrase (standard), au début (insistance), ou après modal (formel).
\item \textbf{Réponses} : \eng{Certainly. / Of course. / Sure.} (Oui) ; \eng{I'm afraid I can't. / Not at all.} (Non / Pour \eng{mind}).
\end{itemize}
\end{aretenir}

\coursec{Règle 2 : Could you / Can you + base verbale (demandes polies)}

Dans la section précédente, nous avons appris à \emph{proposer} notre aide avec \eng{let me + base verbale} : \eng{Let me help you with the accounts.} « Permettez-moi de vous aider pour les comptes. » Mais dans la vie économique — au marché de Mokolo, à la banque, au bureau — il est encore plus fréquent de \emph{demander} quelque chose à quelqu'un. Le français utilise « peux-tu… ? », « pouvez-vous… ? », « pourriez-vous… ? ». L'anglais possède exactement les mêmes nuances avec \cle{can you} et \cle{could you}. Observons d'abord ces structures en situation réelle.

\courssub{Observation : au guichet de la banque}

\begin{textebox}[At the bank in Yaoundé — a dialogue]
\textbf{Customer:} Good morning. Could you tell me my account balance, please?

\textbf{Clerk:} Certainly, sir. Could you give me your account number?

\textbf{Customer:} Of course. It's 402-7781.

\textbf{Clerk:} Thank you. Could you please wait a moment? … Your balance is 85,000 FCFA.

\textbf{Customer:} Thank you. Could you print a statement for me, please?

\textbf{Clerk:} No problem. Can you sign here, please?

\textbf{Customer:} Sure. And can you also change this 10,000 note into smaller notes?

\textbf{Clerk:} I'm afraid I can't, sir. We have no small notes this morning. I'm sorry.

\textbf{Customer:} That's all right. Thank you very much.

\textbf{Clerk:} You're welcome. Have a nice day!
\end{textebox}

\textbf{Glossaire :}

\begin{center}
\begin{tabularx}{\linewidth}{|X|X|}\hline
\rowcolor{popblueL} English & Français \\ \hline
account balance & solde du compte \\ \hline
account number & numéro de compte \\ \hline
statement & relevé de compte \\ \hline
note (banknote) & billet de banque \\ \hline
to change (money) & changer, faire la monnaie \\ \hline
I'm afraid I can't & je crains de ne pas pouvoir (refus poli) \\ \hline
\end{tabularx}
\end{center}

\textbf{Questions d'observation :}
\begin{enumerate}
\item Souligne toutes les demandes polies du dialogue. Quel mot commence chaque question ?
\item Le verbe après \eng{could you} / \eng{can you} porte-t-il un \eng{-s} ou un \eng{to} ?
\item Où est placé \eng{please} dans les phrases ?
\item Comment l'employé accepte-t-il ? Comment refuse-t-il poliment ?
\end{enumerate}

\courssub{La structure : Can you / Could you + base verbale ?}

\begin{propriete}[Forme de la demande polie]
\textbf{Can you / Could you + verbe à l'infinitif sans \eng{to} (base verbale) + complément ?}

\begin{center}
\begin{tabular}{|l|l|l|l|l|}\hline
\rowcolor{popblueL} Auxiliaire + Sujet & \textit{please} (optionnel) & Base verbale & Complément & Ponctuation \\ \hline
Can you & & lend & me 2,000 FCFA & ? \\ \hline
Could you & & sign & this document, please & ? \\ \hline
Could you & please & repeat & the price & ? \\ \hline
\end{tabular}
\end{center}

Trois points rigoureux :
\begin{itemize}
\item Le sujet est \textbf{toujours} \eng{you} (on s'adresse directement à quelqu'un) ; il est placé \textbf{après} l'auxiliaire : c'est une inversion directe, sans \eng{do}.
\item Le verbe reste à la \textbf{base verbale} (infinitif sans \eng{to}), \textbf{jamais} de \eng{-s}, même à la 3\up{e} personne, jamais de \eng{-ing}.
\item \eng{Could} n'est \textbf{pas un passé} ici : c'est un \textbf{conditionnel de politesse}. \eng{Could you help me?} ne veut pas dire « pouvais-tu m'aider ? » mais « pourrais-tu / pourriez-vous m'aider ? ».
\end{itemize}
\end{propriete}

\begin{exemplebox}[Exemples traduits]
\eng{Could you give me a receipt, please?} « Pourriez-vous me donner un reçu, s'il vous plaît ? »

\eng{Can you wait a minute?} « Peux-tu attendre une minute ? »

\eng{Could you please repeat the price?} « Pourriez-vous répéter le prix, s'il vous plaît ? »

\eng{Can you call the manager?} « Peux-tu appeler le directeur ? »

\eng{Could you open an account for me?} « Pourriez-vous m'ouvrir un compte ? »

\eng{Can you lend me your pen?} « Peux-tu me prêter ton stylo ? »
\end{exemplebox}

\courssub{Can you ou Could you ? L'échelle de politesse}

Les deux structures sont correctes, mais elles n'ont pas le même degré de politesse. \eng{Can you} est la demande polie \textbf{standard}, utilisée entre amis, collègues, camarades de classe ou avec un commerçant que l'on connaît. \eng{Could you} est \textbf{plus polie et plus formelle} : on l'emploie avec un supérieur hiérarchique, un client, une personne âgée ou un inconnu.

\begin{popfigure}
\begin{tikzpicture}[font=\small]
\draw[-{Stealth}, very thick, popblue] (0,0) -- (0,6.4) node[above, align=center, text=popdark]{\textbf{plus poli}};
\fill[popgreen] (0,1) circle (2.5pt);
\node[right, align=left, text=popdark] at (0.3,1) {\eng{Open the window!} \quad (impératif sec = impoli)};
\fill[poporange] (0,3.2) circle (2.5pt);
\node[right, align=left, text=popdark] at (0.3,3.2) {\eng{Can you open the window?} \quad (poli, standard)};
\fill[poppurple] (0,5.4) circle (2.5pt);
\node[right, align=left, text=popdark] at (0.3,5.4) {\eng{Could you open the window, please?} \quad (très poli, formel)};
\end{tikzpicture}
\legende{L'échelle de politesse : de l'ordre impoli à la demande très formelle.}
\end{popfigure}

\begin{methode}[Choisir entre can you et could you]
\begin{enumerate}
\item \textbf{À qui parles-tu ?} Identifie ton interlocuteur.
\item \textbf{Ami, camarade, collègue, pair} → utilise \eng{Can you + base verbale ?} \\ \eng{Can you lend me your calculator?}
\item \textbf{Supérieur, client, inconnu, personne âgée, situation officielle} → utilise \eng{Could you + base verbale + please ?} \\ \eng{Could you check my invoice, please?}
\item \textbf{En cas de doute}, choisis \eng{could you} : on n'est jamais trop poli en anglais.
\end{enumerate}
\end{methode}

\courssub{La place de please}

\begin{propriete}[Position de please]
\eng{Please} a deux positions possibles dans la demande :
\begin{itemize}
\item \textbf{en fin de phrase}, précédé d'une virgule : \eng{Could you help me, please?}
\item \textbf{juste après \eng{you}} : \eng{Could you please help me?}
\end{itemize}
Les deux sont correctes. En revanche, on ne place \textbf{jamais} \eng{please} entre l'auxiliaire et le sujet (\eng{Could please you…} est impossible) ni au tout début de la question dans ce type de demande.
\end{propriete}

\courssub{Répondre à une demande polie}

\begin{propriete}[Acceptation et refus poli]
\textbf{Pour accepter :} \eng{Sure.} / \eng{Of course.} / \eng{Certainly.} / \eng{No problem.} / \eng{With pleasure.}

\textbf{Pour refuser poliment :} \eng{I'm sorry, I can't.} / \eng{I'm afraid I can't, I'm busy.} / \eng{I'd like to, but I can't right now.}

Un refus anglais s'accompagne presque toujours d'une excuse (\eng{I'm sorry}, \eng{I'm afraid}) et souvent d'une raison : refuser sèchement par \eng{No.} est très impoli.
\end{propriete}

\begin{attention}[Ne réponds jamais « Yes, I could » !]
Quand on te demande un service avec \eng{Could you…?}, \eng{could} est une politesse, pas une question sur ta capacité. On ne répond donc \textbf{jamais} \eng{Yes, I could} ni \eng{No, I couldn't}. On répond :
\begin{itemize}
\item \eng{Yes, of course.} / \eng{Sure.} / \eng{Certainly.}
\item \eng{I'm sorry, I can't.}
\end{itemize}
De même, répondre \eng{Yes, I can} seul sonne sec ; préfère \eng{Yes, of course}.
\end{attention}

\courssub{Comparaison avec le français et pièges des francophones}

\begin{center}
\begin{tabularx}{\linewidth}{|X|X|X|}\hline
\rowcolor{popblueL} Français & English & Remarque \\ \hline
Peux-tu m'aider ? & Can you help me? & registre familier/standard \\ \hline
Pourriez-vous m'aider ? & Could you help me? & \eng{could} = conditionnel de politesse \\ \hline
S'il vous plaît & please & en fin de phrase ou après \eng{you} \\ \hline
Pouvez-vous… ? & \textbf{jamais} \eng{Do you can…?} & inversion directe, sans \eng{do} \\ \hline
Oui, bien sûr. & Yes, of course. & jamais \eng{Yes, I could} \\ \hline
Je regrette, je ne peux pas. & I'm sorry, I can't. & refus poli avec excuse \\ \hline
\end{tabularx}
\end{center}

\begin{attention}[Erreurs fréquentes des francophones]
\begin{itemize}
\item \eng{Do you can help me?} est \textbf{faux} : les francophones traduisent « pouvez-vous » mot à mot avec \eng{do}. En anglais, l'auxiliaire modal seul fait l'inversion : \eng{Can you help me?}
\item \eng{Could you speaks slowly?} est \textbf{faux} : après \eng{can/could}, le verbe reste à la base verbale, sans \eng{-s} : \eng{Could you speak slowly?}
\item \eng{Could you to help me?} est \textbf{faux} : pas de \eng{to} après un modal.
\item Ne confonds pas \eng{Could you…?} (demande polie au présent) avec le passé \eng{could} dans \eng{When I was young, I could run fast.} Même mot, deux valeurs différentes.
\end{itemize}
\end{attention}

\begin{savaistu}[La politesse dans les banques britanniques]
Dans les banques, les administrations et les magasins britanniques, les employés utilisent presque toujours \eng{Could you…} avec les clients, jamais \eng{Can you…}, jugé trop direct avec un inconnu. On entend toute la journée : \eng{Could you fill in this form, please?}, \eng{Could you wait over there, please?} Cette hyper-politesse fait partie de la culture du service : au Royaume-Uni, un employé qui dirait \eng{Give me your passport} passerait pour extrêmement rude, alors que la phrase est grammaticalement correcte !
\end{savaistu}

\courssub{Entraînement guidé}

\begin{exoresolu}[Mets la demande au bon degré de politesse]
\textbf{Énoncé.} Transforme chaque ordre en demande polie adaptée à la situation.

1. « Give me a receipt. » — un client parle à un caissier qu'il ne connaît pas.

2. « Lend me 2,000 FCFA. » — tu parles à ton meilleur ami.

3. « Repeat the price. » — un client parle à la directrice d'un supermarché de Douala.

4. « Help me carry these boxes. » — tu parles à un camarade de classe.

\tcblower
\textbf{Solution détaillée.}

1. Interlocuteur inconnu, situation commerciale → \eng{could you} + \eng{please} en fin de phrase : \eng{Could you give me a receipt, please?}

2. Meilleur ami → \eng{can you}, sans \eng{please} obligatoire : \eng{Can you lend me 2,000 FCFA?}

3. Supérieure / personne inconnue → \eng{could you} + \eng{please} après \eng{you} : \eng{Could you please repeat the price?}

4. Camarade (pair) → \eng{can you} : \eng{Can you help me carry these boxes?}

\textbf{Méthode :} identifier l'interlocuteur → choisir \eng{can} ou \eng{could} → inversion directe sans \eng{do} → verbe à la base verbale → placer \eng{please} en fin de phrase ou après \eng{you}.
\end{exoresolu}

\begin{exercice}[À toi de jouer]
1. Traduis en anglais : « Pourriez-vous me donner la monnaie, s'il vous plaît ? »

2. Traduis en anglais : « Peux-tu fermer la porte du magasin ? »

3. Corrige la phrase : \eng{Do you can show me the way to the market?}

4. Un client te demande : \eng{Could you wrap this for me, please?} Tu refuses poliment parce que tu n'as plus de papier. Écris ta réponse.

5. Réponds à \eng{Could you sign here, please?} par une acceptation (deux réponses possibles).
\end{exercice}

\begin{corrige}[Exercice « À toi de jouer »]
1. \eng{Could you give me my change, please?} (ou \eng{Could you please give me my change?})

2. \eng{Can you close the shop door?}

3. \eng{Can you show me the way to the market?} — l'inversion se fait avec le modal seul, jamais avec \eng{do}.

4. \eng{I'm afraid I can't, sir. We have no wrapping paper left. I'm sorry.} — refus poli = excuse + \eng{I can't} + raison.

5. \eng{Yes, of course.} / \eng{Certainly.} — jamais \eng{Yes, I could}.
\end{corrige}

\begin{aretenir}[L'essentiel de la règle 2]
\begin{itemize}
\item Demande polie = \eng{Can you / Could you + base verbale (sans to, sans -s) ?}
\item \eng{Can you} = poli standard (amis, pairs) ; \eng{Could you} = plus formel (supérieur, client, inconnu).
\item \eng{Could} est ici un conditionnel de politesse, pas un passé.
\item \eng{Please} se place en fin de phrase ou juste après \eng{you}.
\item On accepte par \eng{Sure. / Of course. / Certainly.} ; on refuse par \eng{I'm sorry, I can't.} / \eng{I'm afraid I can't.}
\item Jamais \eng{Do you can…?} et jamais \eng{Yes, I could} en réponse.
\end{itemize}
\end{aretenir}

\coursec{Comparaison français/anglais et pièges des francophones}

Nous venons de voir que \eng{could you} et \eng{can you} suivis de la base verbale permettent de formuler des demandes polies, et que \eng{let me} exprime une offre d'aide. Mais attention : le français et l'anglais ne construisent pas ces phrases de la même manière. C'est précisément là que les élèves francophones commettent le plus d'erreurs. Observons d'abord quelques paires de phrases, puis dégageons les différences systématiques.

\begin{exemplebox}[Observation : deux langues, deux logiques]
Comparez attentivement :
\begin{itemize}
\item \eng{Let me help you.} = Laisse-moi t'aider. (pas de \eng{to} en anglais, alors que le français dit « t'\textbf{aider} » avec l'infinitif marqué)
\item \eng{Could you lend me some money?} = Pourriez-vous me prêter de l'argent ? (pas de conditionnel en anglais, alors que le français conjugue « pourr\textbf{-iez} »)
\item \eng{Can you open the door, please?} = Peux-tu ouvrir la porte, s'il te plaît ? (\eng{please} en fin de phrase)
\item \eng{Let's go to the market.} = Allons au marché. (\eng{let's} avec apostrophe, pas de \eng{to})
\end{itemize}
\end{exemplebox}

\courssub{Comparaison 1 : « pourriez-vous » (conditionnel) contre \eng{could you} (forme fixe)}

En français, la politesse s'obtient en \textbf{conjuguant le verbe au conditionnel} : « tu peux » devient « tu pourr\textbf{ais} », « vous pouvez » devient « vous pourr\textbf{iez} ». La terminaison change.

En anglais, rien de tel : \eng{could} est une forme \textbf{invariable}. Elle ne prend jamais de terminaison, et le verbe qui suit reste à la base verbale, sans \eng{to}, sans \eng{-s}, sans \eng{-ing}.

\begin{propriete}[Règle : could est invariable]
\eng{Could} + sujet + \cle{base verbale} (demande polie). La forme \eng{could} ne change jamais, quelle que soit la personne :
\begin{itemize}
\item \eng{Could I use your phone?} = Pourrais-je utiliser ton téléphone ?
\item \eng{Could you help me?} = Pourriez-vous m'aider ?
\item \eng{Could she come earlier?} = Pourrait-elle venir plus tôt ?
\end{itemize}
Il n'existe pas de forme \eng{coulds} ni \eng{coulded}. Le verbe principal qui suit ne porte aucune marque : c'est \eng{could} seul qui porte la politesse.
\end{propriete}

\begin{attention}[Piège des francophones : ne pas « conjuguer » l'anglais]
Le francophone veut traduire la terminaison « -iez » de « pourriez ». Résultat fréquent : \eng{Could you to lend me money?} ou \eng{Could you lends me money?} — les deux sont FAUX. En anglais, après \eng{could}, le verbe est \textbf{nu} : \eng{Could you lend me money?}
\end{attention}

\courssub{Comparaison 2 : « s'il vous plaît » et la place de \eng{please}}

Le français place « s'il vous plaît » assez librement : « S'il vous plaît, ouvrez la porte », « Ouvrez la porte, s'il vous plaît », voire « Ouvrez, s'il vous plaît, la porte ».

L'anglais est plus strict : \eng{please} se place \textbf{en fin de phrase} (position la plus courante) ou \textbf{après le sujet \eng{you}}, mais \cle{jamais au milieu du groupe verbal}, jamais entre le modal et le verbe.

\begin{propriete}[Règle : place de please]
\begin{itemize}
\item Fin de phrase : \eng{Could you close the window, please?} = Pourriez-vous fermer la fenêtre, s'il vous plaît ?
\item Après \eng{you} : \eng{Could you please close the window?} (même sens, ton légèrement plus insistant)
\item INTERDIT : \eng{Could you close please the window?} (entre le verbe et son complément)
\end{itemize}
\end{propriete}

\courssub{Les cinq pièges classiques des francophones}

\begin{attention}[Piège 1 : le \eng{to} après \eng{let me}]
Le français dit « laisse-moi \textbf{t}'aider » : l'infinitif français porte une marque. Le francophone la reporte en anglais.
\begin{itemize}
\item FAUX : \eng{Let me to help you.}
\item CORRECT : \eng{Let me help you.} = Laisse-moi t'aider.
\end{itemize}
Après \eng{let me} (et \eng{let us / let's}), le verbe est toujours à la \cle{base verbale}, sans \eng{to}.
\end{attention}

\begin{attention}[Piège 2 : le \eng{to} après \eng{could you}]
\begin{itemize}
\item FAUX : \eng{Could you to lend me money?}
\item CORRECT : \eng{Could you lend me money?} = Pourriez-vous me prêter de l'argent ?
\end{itemize}
\eng{Could} est un modal : le verbe qui suit est nu, comme après \eng{can}, \eng{must}, \eng{should}.
\end{attention}

\begin{attention}[Piège 3 : \eng{do} + modal, le mélange interdit]
Le français forme la question avec un seul verbe : « Peux-tu... ? ». Le francophone, habitué à \eng{do you...?}, ajoute \eng{do} par réflexe.
\begin{itemize}
\item FAUX : \eng{Do you can help me?}
\item CORRECT : \eng{Can you help me?} = Peux-tu m'aider ?
\end{itemize}
Un modal (\eng{can}, \eng{could}, \eng{must}...) fait l'inversion \textbf{tout seul} : il se place devant le sujet, sans \eng{do}. On ne peut jamais avoir \eng{do} et un modal dans la même proposition.
\end{attention}

\begin{exemplebox}[Exemples traduits : à corriger mentalement]
\begin{itemize}
\item FAUX : \eng{Let me to carry your box.} / CORRECT : \eng{Let me carry your box.} = Laisse-moi porter ton carton.
\item FAUX : \eng{Do you can open the door?} / CORRECT : \eng{Can you open the door?} = Peux-tu ouvrir la porte ?
\item FAUX : \eng{Could you helps me with my bags?} / CORRECT : \eng{Could you help me with my bags?} = Pourriez-vous m'aider avec mes sacs ?
\end{itemize}
\end{exemplebox}

\begin{attention}[Piège 4 : répondre \eng{Yes, I could}]
À la question \eng{Could you help me?}, le francophone traduit mot à mot « oui, je pourrais » : \eng{Yes, I could.} Cette réponse est maladroite : elle sonne comme « oui, je pourrais... mais je ne le ferai pas forcément ». En anglais, on répond à une \textbf{demande} par une formule d'acceptation :
\begin{itemize}
\item \eng{Yes, of course.} = Oui, bien sûr.
\item \eng{Sure.} / \eng{Certainly.} = Bien sûr. / Avec plaisir.
\item \eng{No problem.} = Pas de problème.
\item Pour refuser poliment : \eng{I'm sorry, I can't.} = Je suis désolé, je ne peux pas.
\end{itemize}
Rappel : dans \eng{Could you...?}, \eng{could} est une marque de \textbf{politesse}, pas un vrai passé ; la réponse n'a donc pas à reprendre \eng{could}.
\end{attention}

\begin{attention}[Piège 5 : \eng{let's} sans apostrophe ou avec \eng{to}]
« Allons-y », « faisons » se traduisent par \eng{let's} + base verbale. Deux erreurs guettent :
\begin{itemize}
\item Oublier l'apostrophe : \eng{Lets go.} est FAUX (sans apostrophe, \eng{lets} est la 3\textsuperscript{e} personne de \eng{to let}, « il laisse/loue »). CORRECT : \eng{Let's go.} = Allons-y.
\item Ajouter \eng{to} : \eng{Let's to go.} est FAUX. Après \eng{let's}, base verbale nue : \eng{Let's go to the market.} = Allons au marché.
\end{itemize}
\eng{Let's} = \eng{let us} : l'apostrophe remplace le \eng{u} disparu. Elle est obligatoire.
\end{attention}

\courssub{Tableau récapitulatif des erreurs et corrections}

\begin{center}
\begin{tabularx}{\linewidth}{|X|X|X|}
\hline
\rowcolor{popblueL} Erreur fréquente (FAUX) & Correction (CORRECT) & Traduction \\
\hline
Let me to help you. & Let me help you. & Laisse-moi t'aider. \\
\hline
Could you to lend me money? & Could you lend me money? & Pourriez-vous me prêter de l'argent ? \\
\hline
Do you can help me? & Can you help me? & Peux-tu m'aider ? \\
\hline
Could you helps me? & Could you help me? & Pourriez-vous m'aider ? \\
\hline
Yes, I could. (réponse) & Yes, of course. / Sure. & Oui, bien sûr. \\
\hline
Lets go. / Let's to go. & Let's go. & Allons-y. \\
\hline
Could you open please the door? & Could you open the door, please? & Pourriez-vous ouvrir la porte, s'il vous plaît ? \\
\hline
\end{tabularx}
\end{center}

\courssub{Figure récapitulative : correspondances français/anglais}

\begin{popfigure}
\begin{tikzpicture}
\node[draw=popblue, fill=popblueL, rounded corners, minimum width=5.2cm, minimum height=0.9cm, font=\bfseries] (ft) at (0,0) {Français};
\node[draw=popgreen, fill=popgreenL, rounded corners, minimum width=5.2cm, minimum height=0.9cm, font=\bfseries] (et) at (7,0) {Anglais};
\node[draw=popline, fill=popcream, rounded corners, minimum width=5.2cm, minimum height=0.8cm] (f1) at (0,-1.2) {Laisse-moi t'aider.};
\node[draw=popline, fill=popcream, rounded corners, minimum width=5.2cm, minimum height=0.8cm] (e1) at (7,-1.2) {Let me help you.};
\node[draw=popline, fill=popcream, rounded corners, minimum width=5.2cm, minimum height=0.8cm] (f2) at (0,-2.3) {Pourriez-vous m'aider ?};
\node[draw=popline, fill=popcream, rounded corners, minimum width=5.2cm, minimum height=0.8cm] (e2) at (7,-2.3) {Could you help me?};
\node[draw=popline, fill=popcream, rounded corners, minimum width=5.2cm, minimum height=0.8cm] (f3) at (0,-3.4) {S'il vous plaît};
\node[draw=popline, fill=popcream, rounded corners, minimum width=5.2cm, minimum height=0.8cm] (e3) at (7,-3.4) {please (fin de phrase)};
\node[draw=popline, fill=popcream, rounded corners, minimum width=5.2cm, minimum height=0.8cm] (f4) at (0,-4.5) {Allons-y !};
\node[draw=popline, fill=popcream, rounded corners, minimum width=5.2cm, minimum height=0.8cm] (e4) at (7,-4.5) {Let's go!};
\draw[-{Stealth}, thick, poporange] (f1) -- (e1);
\draw[-{Stealth}, thick, poporange] (f2) -- (e2);
\draw[-{Stealth}, thick, poporange] (f3) -- (e3);
\draw[-{Stealth}, thick, poporange] (f4) -- (e4);
\end{tikzpicture}
\legende{Correspondances français/anglais : la politesse se conjugue en français, elle reste fixe en anglais.}
\end{popfigure}

\begin{attention}[La règle d'or : le verbe ne change JAMAIS]
En anglais, le verbe après \eng{let me}, \eng{let's}, \eng{could you} et \eng{can you} ne change \textbf{jamais} : pas de \eng{to}, pas de \eng{-s}, pas de \eng{-ing}. C'est la différence fondamentale avec le français, qui conjugue et marque l'infinitif.
\end{attention}

\begin{methode}[Vérifier sa phrase en trois secondes]
Avant de rendre ta copie ou de parler, applique ce contrôle :
\begin{enumerate}
\item Repère \eng{let me}, \eng{let's}, \eng{could you} ou \eng{can you} dans ta phrase.
\item Regarde le mot qui suit : si tu vois \eng{to}, une terminaison \eng{-s} ou \eng{-ing}, ta phrase est \textbf{fausse}.
\item Supprime la marque fautive : garde uniquement la base verbale (\eng{help}, \eng{lend}, \eng{go}...).
\item Bonus : vérifie l'apostrophe de \eng{let's} et la place de \eng{please} (fin de phrase ou après \eng{you}).
\end{enumerate}
\end{methode}

\begin{exoresolu}[Corrige les erreurs d'un élève]
Énoncé : voici quatre phrases écrites par un élève de Première. Trouve l'erreur dans chacune et réécris la phrase correctement.
\begin{enumerate}
\item \eng{Let me to carry your suitcase, Madam.}
\item \eng{Do you can show me the way to the bank?}
\item \eng{Could you to speak more slowly, please?}
\item \eng{Lets to visit the new market in Buea.}
\end{enumerate}
\tcblower
Solution détaillée :
\begin{enumerate}
\item Après \eng{let me}, pas de \eng{to}. Correction : \eng{Let me carry your suitcase, Madam.} = Laissez-moi porter votre valise, Madame.
\item Un modal ne se combine jamais avec \eng{do} ; \eng{can} fait l'inversion seul. Correction : \eng{Can you show me the way to the bank?} = Pouvez-vous m'indiquer le chemin de la banque ?
\item Après le modal \eng{could}, base verbale nue. Correction : \eng{Could you speak more slowly, please?} = Pourriez-vous parler plus lentement, s'il vous plaît ?
\item Deux erreurs : l'apostrophe manque (\eng{let's}) et \eng{to} est interdit après \eng{let's}. Correction : \eng{Let's visit the new market in Buea.} = Visitons le nouveau marché de Buea.
\end{enumerate}
\end{exoresolu}

\begin{savaistu}[Au baccalauréat : un piège qui rapporte des points]
Au Bac, les exercices de transformation et de réécriture testent très souvent la suppression du \eng{to} après les modaux (\eng{can}, \eng{could}, \eng{must}) et après \eng{let}. Par exemple : « Rewrite the sentence using \eng{could} : \eng{I want you to open the window.} » La réponse attendue est \eng{Could you open the window?} — sans \eng{to}. Les correcteurs repèrent immédiatement les candidats qui maîtrisent la base verbale nue : c'est un critère de correction grammatical qui fait la différence entre une copie moyenne et une bonne copie.
\end{savaistu}

\begin{exercice}[À toi de jouer]
Traduis en anglais correct, en faisant attention aux pièges étudiés :
\begin{enumerate}
\item Laisse-moi t'aider avec tes achats.
\item Pourriez-vous me prêter votre stylo, s'il vous plaît ?
\item Peux-tu fermer la porte ?
\item Allons à la banque !
\item Pourriez-vous répéter, s'il vous plaît ?
\end{enumerate}
\end{exercice}

\begin{corrige}[Exercice « À toi de jouer »]
\begin{enumerate}
\item \eng{Let me help you with your shopping.} (pas de \eng{to} après \eng{let me})
\item \eng{Could you lend me your pen, please?} (\eng{please} en fin de phrase, base verbale \eng{lend})
\item \eng{Can you close the door?} (inversion du modal seul, sans \eng{do})
\item \eng{Let's go to the bank!} (apostrophe obligatoire, pas de \eng{to} après \eng{let's})
\item \eng{Could you repeat, please?} ou \eng{Could you please repeat?} (les deux positions de \eng{please} sont acceptées)
\end{enumerate}
\end{corrige}

\begin{aretenir}[L'essentiel de la section]
\begin{itemize}
\item Le français conjugue la politesse (« pourr\textbf{-iez} ») ; l'anglais utilise des formes fixes : \eng{could you}, \eng{can you}, \eng{let me}, \eng{let's}.
\item Après ces formes, le verbe est toujours à la \cle{base verbale} : jamais de \eng{to}, de \eng{-s} ni de \eng{-ing}.
\item \eng{please} se place en fin de phrase ou après \eng{you}, jamais au milieu du groupe verbal.
\item On répond à une demande par \eng{Yes, of course.} ou \eng{Sure.}, pas par \eng{Yes, I could.}
\item \eng{Let's} prend toujours une apostrophe ; sans elle, le mot change de sens.
\end{itemize}
\end{aretenir}

\coursec{Entraînement guidé et usage en contexte économique}

Maintenant que nous avons comparé le français et l'anglais et repéré les pièges les plus fréquents des francophones (le \eng{to} qui s'invite après \eng{let me}, l'ordre des mots dans \eng{could you}, le \eng{please} mal placé), il est temps de passer à la pratique. Cette dernière étape est essentielle : une règle de grammaire n'est vraiment acquise que lorsqu'elle est utilisée dans des situations réelles. Nous allons d'abord nous entraîner de manière guidée, exercice par exercice, puis utiliser nos structures dans des contextes économiques concrets : le marché, la banque, le bureau, le magasin.

\courssub{Méthode : transformer un ordre en demande polie}

Avant de commencer les exercices, fixons une démarche simple et infaillible, que tu pourras réutiliser au Bac blanc comme à l'examen officiel.

\begin{methode}[Transformer un ordre en demande polie en 4 étapes]
\begin{enumerate}
\item \textbf{Repérer le verbe} de l'ordre (l'impératif) : \eng{Give me your pen.} → verbe : \eng{give}.
\item \textbf{Ajouter \eng{Could you} devant} (ou \eng{Can you} si l'on s'adresse à un pair) : \eng{Could you give...}.
\item \textbf{Garder le verbe sans \eng{to}} et sans changement de forme : jamais \eng{Could you to give} ni \eng{Could you gives}.
\item \textbf{Ajouter \eng{please}} à la fin (ou juste après \eng{you}) : \eng{Could you give me your pen, please?}
\end{enumerate}
\end{methode}

\begin{exemplebox}[Ordre et demande polie]
\begin{itemize}
\item Ordre : \eng{Give me your pen.} → Demande polie : \eng{Could you give me your pen, please?} = « Pourriez-vous me donner votre stylo, s'il vous plaît ? »
\item Ordre : \eng{Wait outside.} → Demande polie : \eng{Could you please wait outside?} = « Pourriez-vous attendre dehors, s'il vous plaît ? »
\item Offre d'aide : \eng{Let me check your bill.} = « Laissez-moi vérifier votre facture. »
\end{itemize}
\end{exemplebox}

\courssub{Exercice 1 : compléter avec let me, let's, could you ou can you}

\begin{exoresolu}[Exercice 1 — À trous]
Énoncé : complète chaque phrase avec \eng{let me}, \eng{let's}, \eng{could you} ou \eng{can you}.

1. \eng{........... carry your bag, Madam. It looks heavy.}

2. \eng{........... go to the market together this afternoon.}

3. \eng{........... lend me your calculator, please? I have a maths test.}

4. \eng{........... help me fill in this form, please? I don't understand question 3.}

5. \eng{........... check the price of the rice on the computer. Wait a moment.}
\tcblower
Solution détaillée :

1. \eng{Let me carry your bag, Madam.} → On propose son aide à quelqu'un : \eng{let me} + base verbale. (« Laissez-moi porter votre sac, Madame. »)

2. \eng{Let's go to the market together this afternoon.} → Suggestion qui inclut le locuteur et les autres : \eng{let's} (= \eng{let us}). (« Allons au marché ensemble cet après-midi. »)

3. \eng{Can you lend me your calculator, please?} → Demande entre camarades (pairs) : \eng{can you} suffit. (« Peux-tu me prêter ta calculatrice, s'il te plaît ? »)

4. \eng{Could you help me fill in this form, please?} → Demande polie à une personne qu'on ne connaît pas bien ou à un supérieur : \eng{could you} est plus poli. (« Pourriez-vous m'aider à remplir ce formulaire, s'il vous plaît ? »)

5. \eng{Let me check the price of the rice on the computer.} → Le vendeur propose de faire quelque chose pour le client : \eng{let me}. (« Laissez-moi vérifier le prix du riz sur l'ordinateur. »)
\end{exoresolu}

\courssub{Exercice 2 : transformer des ordres en demandes polies}

\begin{exercice}[Exercice 2 — De l'ordre à la politesse]
Transforme chaque ordre en demande polie avec \eng{could you} et \eng{please}.

1. \eng{Open the door.}

2. \eng{Show me your identity card.}

3. \eng{Sign here.}

4. \eng{Give me a receipt.}

5. \eng{Speak slowly.}
\end{exercice}

\begin{corrige}[Exercice 2]
1. \eng{Could you open the door, please?} (« Pourriez-vous ouvrir la porte, s'il vous plaît ? »)

2. \eng{Could you show me your identity card, please?} (« Pourriez-vous me montrer votre carte d'identité, s'il vous plaît ? »)

3. \eng{Could you sign here, please?} (« Pourriez-vous signer ici, s'il vous plaît ? »)

4. \eng{Could you give me a receipt, please?} (« Pourriez-vous me donner un reçu, s'il vous plaît ? »)

5. \eng{Could you speak slowly, please?} (« Pourriez-vous parler lentement, s'il vous plaît ? »)

Remarque : le verbe reste toujours à la base verbale, sans \eng{to}, sans \eng{-s}, sans \eng{-ing}.
\end{corrige}

\courssub{Exercice 3 : corriger les erreurs d'élèves francophones}

\begin{exoresolu}[Exercice 3 — Chasse aux erreurs]
Énoncé : chaque phrase contient une erreur typique d'un élève francophone. Corrige-la.

1. \eng{Let me to help you with your luggage.}

2. \eng{Could you to open an account for me?}

3. \eng{You could give me the bill, please?}

4. \eng{Let me explain you the prices.}

5. \eng{Could you repeat please what you said?}
\tcblower
Solution détaillée :

1. \eng{Let me help you with your luggage.} → Pas de \eng{to} après \eng{let me} : \eng{let me} + base verbale. (Le français « laisse-moi \emph{t'}aider » pousse à ajouter \eng{to}.)

2. \eng{Could you open an account for me?} → Pas de \eng{to} après \eng{could you} non plus. (« Pourriez-vous m'ouvrir un compte ? »)

3. \eng{Could you give me the bill, please?} → Inversion obligatoire : auxiliaire \eng{could} + sujet \eng{you}. On ne peut pas garder l'ordre de la phrase affirmative avec une simple intonation comme en français (« Tu pourrais me donner la facture ? »).

4. \eng{Let me explain the prices to you.} → Le verbe \eng{explain} ne prend jamais de complément de personne direct : on dit \eng{explain something to somebody}, jamais \eng{explain somebody}.

5. \eng{Could you repeat what you said, please?} → \eng{Please} se place en fin de phrase (ou juste après \eng{you}), pas au milieu du groupe verbal.
\end{exoresolu}

\begin{attention}[Les 4 erreurs à bannir]
\begin{itemize}
\item \eng{Let me to help you.} → \eng{Let me help you.} (jamais de \eng{to} après \eng{let me} ou \eng{let's}).
\item \eng{Could you to wait?} → \eng{Could you wait?} (base verbale après \eng{could}).
\item \eng{You could help me?} → \eng{Could you help me?} (inversion obligatoire à l'interrogatif).
\item \eng{Please you come here.} → \eng{Could you come here, please?} (\eng{please} ne remplace pas l'auxiliaire ; il ne fait qu'adoucir la demande).
\end{itemize}
\end{attention}

\courssub{Exercice 4 : répondre à des demandes polies}

Savoir demander, c'est bien ; savoir répondre, c'est indispensable. En anglais, on ne répond presque jamais par un simple \eng{yes} : on utilise des formules d'acceptation ou de refus poli.

\begin{center}
\begin{tabularx}{\linewidth}{|X|X|X|}
\hline
\rowcolor{popblueL} Situation & English & Français \\
\hline
Accepter (pair) & \eng{Sure! / Of course! / No problem.} & Bien sûr ! / Avec plaisir ! / Pas de problème. \\
\hline
Accepter (formel) & \eng{Certainly, Madam. / With pleasure.} & Certainement, Madame. / Avec plaisir. \\
\hline
Refuser poliment & \eng{I'm sorry, I can't. I'm busy.} & Je suis désolé, je ne peux pas. Je suis occupé. \\
\hline
Refuser + excuse & \eng{I'm afraid I can't. I have a meeting.} & Je crains de ne pas pouvoir. J'ai une réunion. \\
\hline
\end{tabularx}
\end{center}

\begin{exercice}[Exercice 4 — Accepter ou refuser poliment]
Réponds à chaque demande par une acceptation ou un refus poli, selon l'indication entre parenthèses.

1. \eng{Could you open the window, please?} (accepte)

2. \eng{Can you lend me 5,000 francs?} (refuse poliment)

3. \eng{Could you check my bill, please?} (accepte, formel)

4. \eng{Let me help you with those boxes.} (accepte avec gratitude)

5. \eng{Could you work on Saturday?} (refuse avec une excuse)
\end{exercice}

\begin{corrige}[Exercice 4]
Réponses possibles (d'autres formulations correctes sont acceptées) :

1. \eng{Of course! It's very hot in here.}

2. \eng{I'm sorry, I can't. I don't have any money with me today.}

3. \eng{Certainly, Madam. Let me have a look.}

4. \eng{Thank you very much! That's very kind of you.}

5. \eng{I'm afraid I can't. I have to visit my grandmother in the village.}
\end{corrige}

\begin{attention}[Ne réponds pas par « yes » tout seul]
Dans les dialogues, varie les réponses : \eng{Sure}, \eng{Of course}, \eng{Certainly}, \eng{I'm sorry, I can't}. Un simple \eng{yes} paraît sec, voire impoli en anglais, et te fait perdre des points d'expression orale. Ajoute souvent une petite phrase qui prolonge naturellement l'échange.
\end{attention}

\courssub{Exercice 5 : réécriture type Bac}

\begin{exoresolu}[Exercice 5 — Réécriture avec début imposé]
Énoncé : réécris chaque phrase en commençant par les mots imposés, sans changer le sens.

1. \eng{Help me with this file.} → \eng{Could you ...........................................?}

2. \eng{Carry this bag for me.} → \eng{Can you ...........................................?}

3. \eng{I want to help you fill in the form.} → \eng{Let me ........................................... .}

4. \eng{Shall we start the meeting now?} → \eng{Let's ........................................... .}
\tcblower
Solution détaillée :

1. \eng{Could you help me with this file, please?} → L'ordre devient une demande polie : \eng{could you} + base verbale + \eng{please}.

2. \eng{Can you carry this bag for me, please?} → Demande à un pair : \eng{can you} + base verbale.

3. \eng{Let me help you fill in the form.} → L'offre d'aide se construit avec \eng{let me} + base verbale ; attention, \eng{help you fill in} : après \eng{help}, le deuxième verbe est aussi à la base verbale (sans \eng{to}).

4. \eng{Let's start the meeting now.} → \eng{Shall we...?} et \eng{Let's...} expriment tous deux une suggestion collective.
\end{exoresolu}

\begin{savaistu}[La politesse, clé de la section « Grammar in context »]
Au Bac blanc camerounais comme à l'examen officiel, la section \eng{Grammar in context} propose souvent des dialogues à compléter : un client dans une banque, un passager à l'agence de voyage, un élève chez le proviseur. Dans la plupart des cas, la bonne réponse est la forme \emph{polie} : \eng{Could you...?} plutôt qu'un impératif sec. Entraîne-toi à repérer qui parle à qui : entre amis, \eng{can you} ; à un supérieur ou à un inconnu, \eng{could you... please}. C'est souvent ce choix de registre qui est évalué, autant que la grammaire elle-même.
\end{savaistu}

\courssub{Choisir la bonne formule : organigramme de décision}

Avant de parler, pose-toi deux questions : est-ce que je propose quelque chose, ou est-ce que je demande quelque chose ? Et si je demande, à qui est-ce que je parle ? L'organigramme suivant résume la stratégie.

\begin{popfigure}
\begin{center}
\begin{tikzpicture}[
  node distance=0.9cm and 1.4cm,
  q/.style={rectangle, rounded corners, draw=popblue, fill=popblueL, align=center, font=\small, inner sep=5pt},
  a/.style={rectangle, rounded corners, draw=popgreen, fill=popgreenL, align=center, font=\small, inner sep=5pt},
  lab/.style={font=\scriptsize\itshape, text=popdark},
  arr/.style={-{Stealth[length=2.5mm]}, thick, popdark}
]
\node[q, align=center] (start){Que veux-tu\\exprimer ?};
\node[q, below left=of start, xshift=-1.2cm, align=center] (help){Je veux aider\\ou proposer};
\node[q, below right=of start, xshift=1.2cm, align=center] (ask){Je veux demander\\quelque chose};
\node[a, below=of help, align=center] (letme){\eng{Let me + V}\\\eng{Let's + V} (ensemble)};
\node[q, below=of ask, align=center] (who){À qui est-ce que\\je parle ?};
\node[a, below left=of who, xshift=-0.6cm, align=center] (can){Un pair, un ami\\\eng{Can you + V?}};
\node[a, below right=of who, xshift=0.6cm, align=center] (could){Un supérieur, un inconnu\\\eng{Could you + V, please?}};
\draw[arr] (start) -- (help);
\draw[arr] (start) -- (ask);
\draw[arr] (help) -- (letme);
\draw[arr] (ask) -- (who);
\draw[arr] (who) -- node[lab, left]{pair} (can);
\draw[arr] (who) -- node[lab, right]{supérieur} (could);
\end{tikzpicture}
\end{center}
\legende{Organigramme de décision : choisir entre \eng{let me}, \eng{can you} et \eng{could you} selon l'intention et l'interlocuteur.}
\end{popfigure}

\begin{aretenir}[Tableau récapitulatif des structures]
\begin{center}
\begin{tabularx}{\linewidth}{|l|X|X|}
\hline
\rowcolor{popblueL} Structure & Emploi & Exemple et traduction \\
\hline
\eng{Let me + V} & Offrir son aide & \eng{Let me carry your basket.} « Laisse-moi porter ton panier. » \\
\hline
\eng{Let's + V} & Suggestion collective & \eng{Let's go to the bank.} « Allons à la banque. » \\
\hline
\eng{Can you + V?} & Demande à un pair & \eng{Can you lend me your pen?} « Peux-tu me prêter ton stylo ? » \\
\hline
\eng{Could you + V (+ please)?} & Demande polie, formelle & \eng{Could you wait a moment, please?} « Pourriez-vous attendre un instant, s'il vous plaît ? » \\
\hline
\end{tabularx}
\end{center}
Dans tous les cas, le verbe reste à la \cle{base verbale} : sans \eng{to}, sans \eng{-s}, sans \eng{-ing}.
\end{aretenir}

\courssub{Mise en situation : mini-dialogues en paires}

\begin{experience}[Jeux de rôles : au marché, à la banque, au bureau]
Travaillez par paires. Chaque binôme choisit une situation et prépare un dialogue de 6 à 8 répliques, puis le joue devant la classe. Chaque dialogue doit contenir au moins deux demandes polies (\eng{could you} ou \eng{can you}), une offre d'aide (\eng{let me}) et une suggestion (\eng{let's}).

\textbf{Situation A — At Mokolo market.} Un client veut acheter des tomates et du poisson ; le vendeur propose son aide, le client demande le prix et un sac.

\textbf{Situation B — At the bank in Yaoundé.} Un client veut ouvrir un compte ; l'employé demande la carte d'identité, propose de remplir le formulaire et suggère de s'asseoir.

\textbf{Situation C — At the office.} Un employé demande à son chef la permission de partir plus tôt ; le chef demande de finir un rapport d'abord ; un collègue propose son aide.
\end{experience}

\begin{exemplebox}[Modèle de dialogue — At the bank]
\eng{Customer: Good morning. Could you help me, please? I want to open an account.}

\eng{Clerk: Certainly, Sir. Could you show me your identity card, please?}

\eng{Customer: Of course. Here it is.}

\eng{Clerk: Thank you. Let me give you the application form.}

\eng{Customer: Could you help me fill it in, please? My English is not very good.}

\eng{Clerk: No problem. Let's sit down at this desk.}

\eng{Customer: Thank you very much.}

\eng{Clerk: You're welcome, Sir.}
\end{exemplebox}

\courssub{Production guidée : une journée de travail dans un magasin}

\begin{exercice}[Exercice 6 — Production écrite guidée]
Tu travailles pour la journée dans un magasin de vêtements à Douala. Rédige cinq demandes polies que tu pourrais adresser aux clients ou à ton employeur, en utilisant au moins une fois chacune des structures \eng{could you}, \eng{can you}, \eng{let me} et \eng{let's}. Souligne la structure utilisée dans chaque phrase.
\end{exercice}

\begin{corrige}[Exercice 6 — Exemple de production]
1. \eng{Could you wait a moment, please? The shop is very busy.} (à un client)

2. \eng{Let me show you our new shirts, Madam.} (offre d'aide à une cliente)

3. \eng{Could you sign this delivery note, please?} (au livreur)

4. \eng{Can you count the money in the till, please?} (à un collègue)

5. \eng{Let's close the shop now; it is six o'clock.} (suggestion aux collègues)

Critères de réussite : la base verbale après chaque structure, l'inversion \eng{could you} / \eng{can you} correcte, \eng{please} bien placé, et un vocabulaire adapté au contexte du magasin (\eng{till}, \eng{delivery note}, \eng{customer}).
\end{corrige}

\begin{aretenir}[L'essentiel de la leçon]
Pour aider ou proposer : \eng{let me + V} (aide) et \eng{let's + V} (suggestion collective). Pour demander : \eng{can you + V?} entre pairs, \eng{could you + V, please?} pour plus de politesse. Le verbe reste toujours à la base verbale, l'inversion est obligatoire dans les questions, et les réponses se varient : \eng{Sure}, \eng{Of course}, \eng{Certainly}, \eng{I'm sorry, I can't}. Ces formules, simples et courtes, sont parmi les plus rentables de tout le programme de Première : elles servent à l'oral, à l'écrit et dans la section \eng{Grammar in context} de l'examen.
\end{aretenir}

\newpage

\coursec{Activité d'intégration}

\coursec{Lesson 14 — Grammar: Polite requests and suggestions}

\courssub{Objectifs d'apprentissage}

À l'issue de cette leçon, tu seras capable de :

\begin{itemize}
\item comprendre un dialogue de négociation commerciale authentique et identifier les formules de politesse ;
\item construire correctement les structures \cle{let me} / \cle{let's} + \cle{base verbale} pour proposer ton aide ou suggérer une action commune ;
\item formuler des demandes polies avec \cle{could you} (registre soutenu) et \cle{can you} (registre courant), suivies de la base verbale ;
\item placer correctement \cle{please} (début ou fin de phrase, jamais devant le verbe) ;
\item accepter (\eng{Of course, With pleasure}) ou refuser poliment (\eng{I'm sorry, I can't, but...}) une proposition ;
\item produire oralement puis par écrit un dialogue de 8 à 10 répliques dans un contexte économique camerounais réaliste (achat, prix, négociation).
\end{itemize}

\courssub{Situation de départ : Au marché de Buea}

\begin{textebox}[At the Buea Market — Real-life situation]
It is Saturday morning at the Buea main market. The stalls are full of shoes, bags, second-hand clothes and fresh food. Mr Tabi, a shoe seller, has just arranged his best pairs of trainers and school shoes on a wooden table. A young customer, Miss Ngum, stops in front of his stall. She needs a good pair of black shoes for her new job at a microfinance office in town, but her budget is limited: she has only 15 000 francs CFA. The shoes she likes cost 18 000 francs. She wants to negotiate politely. Mr Tabi wants to help her, offer his products, and make a good sale without losing too much money. Both of them must be polite, because politeness is the key to a successful negotiation in Cameroon.
\end{textebox}

\begin{definition}[Glossaire utile — Vocabulaire de l'achat et de la négociation]
\begin{itemize}
\item \eng{stall} (n.) : étal, stand de marché. \quad \eng{trainer} (n.) : chaussure de sport, basket (UK).
\item \eng{budget} (n.) : budget, somme disponible. \quad \eng{to negotiate} (v.) : négocier.
\item \eng{a discount} (n.) : une remise, une réduction. \quad \eng{a bargain} (n.) : une bonne affaire.
\item \eng{to afford} (v.) : avoir les moyens de s'acheter. \quad \eng{customer} (n.) : client(e).
\item \eng{price} (n.) : prix. \quad \eng{cheap / expensive} (adj.) : bon marché / cher.
\item \eng{cash} (n.) : espèces, argent liquide. \quad \eng{change} (n.) : monnaie (rendue).
\item \eng{size} (n.) : pointure, taille. \quad \eng{quality} (n.) : qualité.
\item \eng{to wrap} (v.) : emballer. \quad \eng{receipt} (n.) : reçu.
\end{itemize}
\end{definition}

\courssub{Étape 1 : Compréhension et observation (Travail en binôme)}

\begin{enumerate}
\item \textbf{Lecture à deux voix.} Lisez le texte ci-dessus. Répondez en anglais :
      \eng{Who are the two speakers? What does the customer want? How much money does she have? What is the problem?}
\item \textbf{Recherche lexicale.} Avec le glossaire, listez 8 mots anglais indispensables pour parler d'achat et de prix (ex. \eng{price, discount, cash...}).
\item \textbf{Observation guidée.} Relisez le dialogue modèle (page suivante) et soulignez :
      \begin{itemize}
      \item deux offres d'aide avec \eng{let me} ;
      \item une suggestion commune avec \eng{let's} ;
      \item deux demandes très polies avec \eng{could you ... please} ;
      \item une demande plus directe avec \eng{can you} ;
      \item une acceptation et un refus poli.
      \end{itemize}
\end{enumerate}

\courssub{Règle 1 : \texorpdfstring{\eng{Let me} / \eng{Let's} + base verbale}{Let me / Let's + bare infinitive} — Offres d'aide et suggestions}

\begin{propriete}[Forme et emploi]
\begin{tabular}{|l|l|l|}
\hline
\rowcolor{popblueL} \textbf{Structure} & \textbf{Emploi} & \textbf{Exemple} \\
\hline
\eng{Let me} + \cle{base verbale} & Offre d'aide personnelle (le locuteur agit). & \eng{Let me help you.} (« Permettez-moi de vous aider. ») \\
\hline
\eng{Let's} (= \eng{Let us}) + \cle{base verbale} & Suggestion d'action commune (locuteur + interlocuteur). & \eng{Let's discuss the price.} (« Discutons le prix. ») \\
\hline
\end{tabular}

\medskip
\textbf{Règles de formation :}
\begin{itemize}
\item \cle{Pas de \eng{to}} après \eng{let me} / \eng{let's}. \eng{Let me show you} (correct) — \eng{Let me to show you} (incorrect).
\item Le verbe reste à la \cle{base verbale} (infinitif sans \eng{to}) : \eng{go, see, make, wrap, say}.
\item \eng{Let's} est la contraction de \eng{let us} ; à l'oral : « lets » (une syllabe, \eng{/lets/}).
\item Négation : \eng{Let's not} + base verbale (ex. \eng{Let's not argue.}). Pour \eng{let me}, on n'utilise pas la forme négative pour une offre d'aide.
\end{itemize}
\end{propriete}

\begin{exemplebox}[Exemples variés — Contexte économique]
\begin{itemize}
\item \eng{Let me carry that bag for you.} (Vendeur aide le client)
\item \eng{Let me calculate the total.} (Caissier)
\item \eng{Let's meet tomorrow at 10 a.m. to sign the contract.} (Partenaires)
\item \eng{Let's say 15 000 francs and we have a deal.} (Négociation : proposition de prix)
\item \eng{Let me give you my business card.} (Networking)
\end{itemize}
\end{exemplebox}

\courssub{Règle 2 : \texorpdfstring{\eng{Could you} / \eng{Can you} + base verbale}{Could you / Can you + bare infinitive} — Demandes polies}

\begin{propriete}[Degrés de politesse et forme]
\begin{tabular}{|l|l|l|}
\hline
\rowcolor{popblueL} \textbf{Structure} & \textbf{Registre / Nuance} & \textbf{Exemple} \\
\hline
\eng{Could you} + base verbale (+ \eng{please}) & Très poli, formel, recommandé avec un inconnu ou un supérieur. & \eng{Could you show me the black pair, please?} \\
\hline
\eng{Can you} + base verbale (+ \eng{please}) & Poli mais plus direct, courant entre collègues ou après un premier contact. & \eng{Can you give me a discount?} \\
\hline
\end{tabular}

\medskip
\textbf{Règles de formation :}
\begin{itemize}
\item \cle{Pas de \eng{to}} après \eng{could you} / \eng{can you}. \eng{Could you come?} (correct) — \eng{Could you to come?} (incorrect).
\item Le verbe principal reste à la \cle{base verbale} : \eng{reduce, add, wrap, call, send}.
\item \eng{Could} est le prétérit modal de \eng{can} ; il exprime la distance (politesse), pas le passé.
\item \eng{Please} : position \cle{initiale} (\eng{Please, could you...}) ou \cle{finale} (\eng{...could you... please?}). \textbf{Jamais} devant le verbe (\eng{*Could you please show me} est possible mais moins naturel en UK standard que \eng{Could you show me, please?}).
\end{itemize}
\end{propriete}

\begin{exemplebox}[Demandes types en situation professionnelle]
\begin{itemize}
\item \eng{Could you send me the catalogue by email, please?}
\item \eng{Can you deliver the goods before Friday?}
\item \eng{Could you repeat the price, please? I didn't catch it.}
\item \eng{Can you hold the line a moment?} (Téléphone)
\item \eng{Could you please sign here?} (Position médiane de \eng{please} acceptée dans l'administratif).
\end{itemize}
\end{exemplebox}

\courssub{Réponses types : Accepter et Refuser poliment}

\begin{propriete}[Réagir à une demande / offre / suggestion]
\begin{center}
\begin{tabularx}{\linewidth}{|X|X|}
\hline
\rowcolor{popgreenL} \textbf{Accepter (Accepting)} & \textbf{Refuser poliment (Refusing politely)} \\
\hline
\eng{Of course.} / \eng{Certainly.} & \eng{I'm sorry, I can't.} \\
\eng{With pleasure.} / \eng{My pleasure.} & \eng{I'm afraid that's not possible.} \\
\eng{No problem.} / \eng{Sure.} (informel) & \eng{I'm afraid I can't, but...} (contre-proposition) \\
\eng{Yes, let's do that.} (pour \eng{Let's}) & \eng{That's a bit too high/low for me, but...} \\
\hline
\end{tabularx}
\end{center}
\medskip
\cle{Astuce} : Un refus poli en négociation s'accompagne \textbf{toujours} d'une explication ou d'une contre-proposition (\eng{but...}).
\end{propriete}

\courssub{Comparaison Français / Anglais — Pièges des francophones}

\begin{attention}[Erreurs fréquentes à ne plus commettre]
\begin{itemize}
\item \textbf{Faux ami structurel :} \eng{Let me to help you} $\leftarrow$ \textbf{INCORRECT}. \fr{« Permettez-moi de vous aider »} $\rightarrow$ \eng{Let me help you.} (Pas de \eng{to}).
\item \textbf{Faux ami structurel :} \eng{Could you to come?} $\leftarrow$ \textbf{INCORRECT}. \fr{« Pourriez-vous venir ? »} $\rightarrow$ \eng{Could you come?} (Pas de \eng{to} après modal).
\item \textbf{Traduction mot à mot :} \eng{If you please} ou \eng{Please you} $\leftarrow$ \textbf{INCORRECT}. \fr{« S'il vous plaît »} $\rightarrow$ \eng{Please} (seul, en début ou fin de phrase).
\item \textbf{Registre :} Utiliser \eng{Can you} avec un client important ou un supérieur au premier abord $\leftarrow$ \textbf{MALADROIT}. Préférer \eng{Could you ... please?}.
\item \textbf{Ordre des mots :} \eng{*Could you please to help me?} $\leftarrow$ Double erreur (\eng{to} + place de \eng{please}). Correct : \eng{Could you help me, please?}
\item \textbf{Prononciation :} \eng{Let's} se prononce « lets » (voyelle \eng{/e/} comme dans \eng{bed}, \textbf{pas} « laits »). \eng{Could you} $\rightarrow$ « coud-you » (liaison \eng{/d/} + \eng{/j/} $\rightarrow$ \eng{/kʊdʒuː/}).
\end{itemize}
\end{attention}

\begin{center}
\begin{tabularx}{\linewidth}{|X|X|X|}
\hline
\rowcolor{poporangeL} \textbf{Français} & \textbf{English (Correct)} & \textbf{Piège à éviter} \\
\hline
Permettez-moi de voir. & \eng{Let me see.} & \eng{Let me to see.} \\
Discutons le prix. & \eng{Let's discuss the price.} & \eng{Let's to discuss...} \\
Pourriez-vous baisser le prix ? & \eng{Could you reduce the price, please?} & \eng{Could you to reduce...} \\
Pouvez-vous m'aider ? & \eng{Can you help me?} & \eng{Can you to help me?} \\
S'il vous plaît, montrez-moi. & \eng{Please, show me.} / \eng{Show me, please.} & \eng{Please you show me.} \\
\hline
\end{tabularx}
\end{center}

\courssub{Entraînement guidé : Exercices d'application}

\begin{exercice}[Transformation et construction]
\begin{enumerate}
\item Transforme les phrases suivantes en offres d'aide avec \eng{Let me} :
      \begin{itemize}
      \item I will carry your box. $\rightarrow$ \dots
      \item I can explain the contract. $\rightarrow$ \dots
      \item I'll wrap the gift. $\rightarrow$ \dots
      \end{itemize}
\item Transforme en suggestions communes avec \eng{Let's} :
      \begin{itemize}
      \item We should meet tomorrow. $\rightarrow$ \dots
      \item We can negotiate the salary. $\rightarrow$ \dots
      \item Why don't we order pizza? $\rightarrow$ \dots
      \end{itemize}
\item Mets au niveau de politesse supérieur (utilise \eng{Could you ... please}) :
      \begin{itemize}
      \item Can you open the window? $\rightarrow$ \dots
      \item Give me the bill, please. $\rightarrow$ \dots
      \item Can you wait a minute? $\rightarrow$ \dots
      \end{itemize}
\item Complète avec \eng{let me}, \eng{let's}, \eng{could you} ou \eng{can you} + verbe entre parenthèses :
      \begin{enumerate}
      \item \dots (\eng{show}) you the new model? It just arrived.
      \item \dots (\eng{go}) to the bank first, then to the market.
      \item \dots (\eng{calculate}) the discount for you, Sir?
      \item \dots (\eng{not}) \dots (\eng{waste}) time; the shop closes at 6 p.m.
      \item \dots (\eng{add}) a warranty on this laptop?
      \end{enumerate}
\end{enumerate}
\end{exercice}


\begin{corrige}[Exercice 1 -- Transformation et construction]
\begin{enumerate}
\item \eng{Let me carry your box.} / \eng{Let me explain the contract.} / \eng{Let me wrap the gift.}
\item \eng{Let's meet tomorrow.} / \eng{Let's negotiate the salary.} / \eng{Let's order pizza.}
\item \eng{Could you open the window, please?} / \eng{Could you give me the bill, please?} / \eng{Could you wait a minute, please?}
\item \begin{enumerate}
      \item \eng{Could you show} (demande polie client/vendeur) ou \eng{Let me show} (offre vendeur).
      \item \eng{Let's go} (suggestion commune).
      \item \eng{Let me calculate} (offre vendeur) ou \eng{Could you calculate} (demande client).
      \item \eng{Let's not waste} (suggestion négative).
      \item \eng{Could you add} (demande polie) ou \eng{Can you add} (plus direct).
      \end{enumerate}
\end{enumerate}
\end{corrige}

\courssub{Production orale guidée : Jeu de rôle « Négociation au marché »}

\begin{experience}[Scénario : Marché de Buea / Douala / Yaoundé]
\textbf{Contexte :} Un vendeur (articles : chaussures, téléphones, tissus, cacao) et un client (budget limité, besoin précis).
\textbf{Contraintes linguistiques obligatoires dans le dialogue (8--10 répliques) :}
\begin{itemize}
\item 1 offre d'aide avec \eng{Let me} + BV.
\item 1 suggestion avec \eng{Let's} + BV.
\item 2 demandes avec \eng{Could you ... please?} + BV.
\item 1 demande avec \eng{Can you ...?} + BV.
\item 1 acceptation polie (\eng{Of course / With pleasure}).
\item 1 refus poli avec contre-proposition (\eng{I'm sorry, I can't, but...}).
\item Vocabulaire : \eng{price, discount, cash, quality, size, wrap, receipt, bargain}.
\end{itemize}
\textbf{Déroulement :}
\begin{enumerate}
\item Préparation 10 min (écriture au brouillon, vérification structures).
\item Répétition 5 min (intonation montante sur les questions \eng{Could/Can you...?}, ton aimable).
\item Présentation devant la classe (les meilleurs binômes).
\item Évaluation par les pairs : grille « Structures correctes / Prononciation / Fluidité / Vocabulaire éco ».
\end{enumerate}
\end{experience}

\courssub{Production écrite modèle et analyse}

\begin{exoresolu}[At the Buea Market — Dialogue modèle annoté]
Énoncé : Rédige un dialogue de 8 à 10 répliques entre Mr Tabi (seller) et Miss Ngum (customer) contenant au moins deux offres/suggestions avec \eng{let me} ou \eng{let's}, deux demandes avec \eng{could you ... please}, une demande avec \eng{can you}, une acceptation et un refus poli avec contre-proposition.
\tcblower
\textbf{Modèle de production :}

\textbf{Mr Tabi:} \eng{Good morning, Madam! Welcome to my stall. Let me help you. Are you looking for anything special?}

\textbf{Miss Ngum:} \eng{Good morning. Yes, I need a pair of black shoes for my new job. Could you show me those leather ones on the table, please?}

\textbf{Mr Tabi:} \eng{Of course! Here you are. They are genuine leather, very durable. Let me tell you, the quality is excellent.}

\textbf{Miss Ngum:} \eng{They look smart. How much do they cost?}

\textbf{Mr Tabi:} \eng{They are 18 000 francs CFA, Madam. A real bargain for this quality.}

\textbf{Miss Ngum:} \eng{Hmm, that is a little expensive for me. I have only 15 000 francs. Could you reduce the price, please?}

\textbf{Mr Tabi:} \eng{I'm sorry, I can't go down to 15 000, but let me make you a special offer. Let's say 16 500 francs.}

\textbf{Miss Ngum:} \eng{Can you add a free pair of socks at that price?}

\textbf{Mr Tabi:} \eng{With pleasure! You are a good negotiator. Let me wrap everything for you.}

\textbf{Miss Ngum:} \eng{Thank you very much. Here is the cash. Let's do business again next time!}

\textbf{Mr Tabi:} \eng{Certainly, Madam. Have a nice day!}

\medskip
\textbf{Analyse linguistique commentée :}
\begin{itemize}
\item \eng{Let me help you / Let me tell you / Let me make you a special offer / Let me wrap everything} : 4 offres d'aide du vendeur (\eng{Let me} + BV). Aucune trace de \eng{to}.
\item \eng{Let's say 16 500 francs / Let's do business again} : 2 suggestions communes (\eng{Let's} + BV), typiques de la négociation gagnant-gagnant.
\item \eng{Could you show me ... please? / Could you reduce the price, please?} : 2 demandes très polies de la cliente (\eng{Could you} + BV + \eng{please} en fin).
\item \eng{Can you add a free pair of socks?} : Demande plus directe (\eng{Can you} + BV), légitime une fois la relation établie.
\item \eng{Of course! / With pleasure! / Certainly.} : Acceptations franches et courtoises.
\item \eng{I'm sorry, I can't go down to 15 000, but let me make you a special offer.} : Refus poli \textbf{suivi d'une contre-proposition} (marqueur de négociation courtoise).
\item \eng{genuine leather, durable, bargain, cash, wrap} : Lexique économique ciblé Chapitre 2.
\end{itemize}
\end{exoresolu}

\courssub{Bilan de la leçon — Ce que tu dois maîtriser}

\begin{aretenir}[Synthèse grammaticale et communicative]
\begin{itemize}
\item \textbf{Offre d'aide :} \eng{Let me} + \cle{base verbale} (sans \eng{to}). Ex. \eng{Let me pay.}
\item \textbf{Suggestion commune :} \eng{Let's} + \cle{base verbale}. Ex. \eng{Let's sign the contract.} Négation : \eng{Let's not argue.}
\item \textbf{Demande très polie :} \eng{Could you} + \cle{base verbale} + \eng{please?} (inconnu, supérieur, client).
\item \textbf{Demande polie directe :} \eng{Can you} + \cle{base verbale} + \eng{please?} (connaissance, collègue).
\item \textbf{Place de \eng{please} :} Début ou fin de phrase. \eng{Please, could you...?} / \eng{Could you... please?}
\item \textbf{Réponses :} Acceptation \eng{Of course / Certainly / With pleasure.} Refus \eng{I'm sorry, I can't, but...}
\item \textbf{Contexte éco (Chap. 2) :} Acheter, vendre, négocier \eng{price, discount, cash, change, quality, size, receipt, bargain}.
\end{itemize}
\end{aretenir}

\begin{savaistu}[Culture de la politesse : Cameroun et monde anglophone]
Au Cameroun anglophone (Northwest, Southwest) comme au Nigeria, au Ghana ou au Royaume-Uni, la \cle{politesse verbale} est un capital social. On ne dit pas « Donne-moi ça » (\eng{Give me that}) mais \eng{Could I have that, please?} ou \eng{Let me have that, please.}
\begin{itemize}
\item \eng{Madam / Sir} s'utilise systématiquement au marché, à la banque, à l'administration pour s'adresser à un adulte.
\item \eng{Please} et \eng{Thank you} (ou \eng{Thanks}) sont obligatoires à chaque tour de parole transactionnel.
\item La négociation (\eng{bargaining / haggling}) au marché est un \textbf{jeu social codifié} : on commence par \eng{Could you reduce...?}, le vendeur refuse poliment (\eng{I'm afraid not... but let's say...}), le client contre-propose (\eng{Can you add...?}). Ne pas négocier peut passer pour de l'arrogance ; trop insister sans \eng{please} passe pour de l'impolitesse.
\item \eng{Let me} est la marque du \textbf{service client} : le bon vendeur dit \eng{Let me check / Let me wrap / Let me carry}.
\end{itemize}
\end{savaistu}

\courssub{Auto-évaluation finale}

Coche les cases que tu maîtrises parfaitement (objectif : tout cocher avant le contrôle continu) :

\begin{tabular}{|l|c|}
\hline
\rowcolor{popblueL} \textbf{Compétence} & \textbf{Acquis} \\
\hline
Je construis \eng{Let me / Let's} + BV sans \eng{to}. & \square \\
\hline
Je forme \eng{Could you / Can you} + BV sans \eng{to}. & \square \\
\hline
Je place \eng{please} correctement (début/fin). & \square \\
\hline
Je choisis \eng{Could you} (formel) vs \eng{Can you} (courant). & \square \\
\hline
J'accepte poliment (\eng{Of course, With pleasure}). & \square \\
\hline
Je refuse poliment avec \eng{I'm sorry, I can't, but...}. & \square \\
\hline
J'utilise le vocabulaire éco : \eng{price, discount, cash, bargain, wrap, receipt}. & \square \\
\hline
Je joue un dialogue de négociation de 8--10 répliques fluide. & \square \\
\hline
\end{tabular}

\newpage

\coursec{Méthodes et exercices résolus}

\courssub{Fiche méthode 1 — Identifier la structure demandée : suggestion, offre d'aide ou demande polie ?}

\begin{methode}[Choisir la bonne structure de politesse]
Avant de construire une phrase, il faut d'abord \emph{identifier la fonction} demandée. Procède en trois étapes :
\begin{enumerate}
\item \textbf{Lis la consigne ou le contexte.} Cherche les mots-clés : \eng{suggest}, \eng{offer to help} indiquent \cle{let me} ; \eng{ask politely}, \eng{request} indiquent \cle{could you} / \cle{can you}.
\item \textbf{Repère le sujet de l'action.} Si \emph{je} fais l'action pour aider quelqu'un, c'est \eng{Let me + base verbale}. Si \emph{tu} demandes à l'autre de faire l'action, c'est \eng{Could you / Can you + base verbale}.
\item \textbf{Choisis le degré de politesse.} \eng{Can you} est courant entre amis ; \eng{Could you} est plus poli et plus formel (clients, supérieurs, examen écrit). Ajoute \eng{please} pour renforcer la politesse.
\end{enumerate}
\textbf{Réflexes à retenir :}
\begin{center}
\begin{tabularx}{\linewidth}{|X|X|X|}\hline
\rowcolor{popblueL} Fonction & Structure & Exemple \\ \hline
Offre d'aide / suggestion & Let me + BV & Let me help you. \\ \hline
Demande polie (formelle) & Could you + BV (please)? & Could you open the door, please? \\ \hline
Demande polie (familière) & Can you + BV? & Can you lend me your pen? \\ \hline
\end{tabularx}
\end{center}
\end{methode}

\begin{exoresolu}[Identifier la fonction et la structure]
\textbf{Énoncé.} For each situation, say which structure is needed: \eng{Let me}, \eng{Could you} or \eng{Can you}. Do not write the full sentence yet.\\
1. You see an old woman carrying heavy bags at the Mokolo market. You want to help her.\\
2. You are in a bank in Yaoundé. You want the clerk to check your account balance.\\
3. You phone your friend. You want him to call you back this evening.\\
4. You are a shopkeeper. A customer looks tired. You propose to carry his goods to his car.
\tcblower
\textbf{Solution détaillée.}
\begin{enumerate}
\item \textbf{Situation 1 :} c'est toi qui fais l'action pour aider quelqu'un $\Rightarrow$ offre d'aide $\Rightarrow$ \textbf{Let me}. (\eng{Let me carry your bags.})
\item \textbf{Situation 2 :} tu demandes à un employé de banque (contexte formel, inconnu) de faire quelque chose $\Rightarrow$ demande polie formelle $\Rightarrow$ \textbf{Could you}. (\eng{Could you check my balance, please?})
\item \textbf{Situation 3 :} tu demandes à un ami (contexte familier) $\Rightarrow$ demande polie familière $\Rightarrow$ \textbf{Can you}. (\eng{Can you call me back this evening?})
\item \textbf{Situation 4 :} le commerçant propose de faire lui-même l'action $\Rightarrow$ offre d'aide $\Rightarrow$ \textbf{Let me}. (\eng{Let me carry your goods to your car.})
\end{enumerate}
\textbf{Conclusion :} la clé est toujours la question « qui fait l'action ? » : moi $\Rightarrow$ \eng{Let me} ; l'autre $\Rightarrow$ \eng{Could you / Can you}. Le degré de formalité tranche ensuite entre \eng{could} (formel) et \eng{can} (familier).
\end{exoresolu}

\courssub{Fiche méthode 2 — Construire correctement \eng{Let me + base verbale}}

\begin{methode}[Construire une offre d'aide ou une suggestion avec \eng{Let me}]
\begin{enumerate}
\item Écris \cle{Let me} (jamais \eng{let I} : \eng{me} est un pronom complément).
\item Ajoute directement la \cle{base verbale} (l'infinitif sans \eng{to}) : \eng{Let me \textbf{help} you.} Pas de \eng{to}, pas de \eng{-ing}, pas de conjugaison.
\item Complète avec le complément : \eng{Let me carry your luggage.} — « Laisse-moi porter tes bagages. »
\item Pour la \textbf{forme négative} : \eng{Let me \textbf{not} + BV} (ex. \eng{Let me not disturb you.}) — attention, on ne dit jamais \eng{don't let me} pour une offre d'aide.
\item Variante pour inclure l'autre personne : \eng{\textbf{Let's} + BV} (= \eng{Let us}) pour une suggestion commune : \eng{Let's go to the market.} — « Allons au marché. »
\end{enumerate}
\textbf{Structure à retenir :} \fbox{Let me + base verbale (+ complément)} / \fbox{Let's + base verbale}
\end{methode}

\begin{exoresolu}[Transformer des phrases avec \eng{Let me}]
\textbf{Énoncé.} Rewrite each sentence using \eng{Let me} or \eng{Let's}, as in the example.\\
Example: I want to pay for your taxi. $\rightarrow$ Let me pay for your taxi.\\
1. I want to show you the way to the post office.\\
2. I propose that we visit the new supermarket in Douala.\\
3. I want to explain the prices to the customers.\\
4. I suggest that we start the meeting now.\\
5. I want to carry the bag of rice for you, madam.
\tcblower
\textbf{Solution détaillée.}
\begin{enumerate}
\item \eng{Let me show you the way to the post office.} — \eng{me} car c'est « je » qui agis seul ; base verbale \eng{show} sans \eng{to}.
\item \eng{Let's visit the new supermarket in Douala.} — « nous » faisons l'action ensemble $\Rightarrow$ \eng{Let's} (= \eng{let us}), base verbale \eng{visit}.
\item \eng{Let me explain the prices to the customers.} — attention à la préposition : \eng{explain something \textbf{to} somebody}.
\item \eng{Let's start the meeting now.} — suggestion collective $\Rightarrow$ \eng{Let's}.
\item \eng{Let me carry the bag of rice for you, madam.} — offre d'aide individuelle ; \eng{carry} reste à la base verbale.
\end{enumerate}
\textbf{Conclusion :} après \eng{let me / let's}, le verbe est toujours à la \emph{base verbale}, sans \eng{to} et sans terminaison. Le choix entre \eng{let me} et \eng{let's} dépend uniquement du nombre de personnes qui font l'action.
\end{exoresolu}

\courssub{Fiche méthode 3 — Construire une demande polie avec \eng{Could you / Can you}}

\begin{methode}[Formuler une demande polie correcte]
\begin{enumerate}
\item Commence par \cle{Could you} (formel, poli) ou \cle{Can you} (familier).
\item Ajoute la \cle{base verbale} directement : \eng{Could you \textbf{repeat}, please?} — jamais \eng{to repeat}, jamais \eng{repeats}.
\item Place \eng{please} en fin de phrase (ou après \eng{you}) : \eng{Could you \textbf{please} help me?} / \eng{Could you help me, \textbf{please}?}
\item Pour une demande encore plus douce, utilise \eng{Could you \textbf{possibly} + BV ?} : \eng{Could you possibly lend me some money?}
\item \textbf{Réponses types :} accepter : \eng{Sure. / Of course. / Certainly. / No problem.} ; refuser poliment : \eng{I'm sorry, I can't. / I'm afraid I can't.}
\end{enumerate}
\textbf{Attention :} \eng{could} ici n'est pas un passé ! C'est un \emph{conditionnel de politesse}. La question porte sur le présent.
\end{methode}

\begin{exoresolu}[Construire des demandes polies en contexte économique]
\textbf{Énoncé.} You are a new employee in a company in Buea. Make polite requests for the following situations, using \eng{Could you} or \eng{Can you} + \eng{please}.\\
1. You want the secretary to give you the price list.\\
2. You want a customer to fill in a form.\\
3. You want your colleague to lend you a calculator.\\
4. You want the manager to sign a document.\\
5. You want a client to wait a moment.
\tcblower
\textbf{Solution détaillée.}
\begin{enumerate}
\item \eng{Could you give me the price list, please?} — contexte formel (secrétaire, nouveau poste) $\Rightarrow$ \eng{could} ; base verbale \eng{give}.
\item \eng{Could you fill in this form, please?} — verbe à particule : \eng{fill in} (= remplir) ; la particule reste accolée au verbe.
\item \eng{Can you lend me a calculator, please?} — entre collègues, \eng{can} est acceptable ; attention : \eng{lend} (prêter à quelqu'un), pas \eng{borrow} (emprunter à quelqu'un).
\item \eng{Could you sign this document, please?} — supérieur hiérarchique $\Rightarrow$ \eng{could} obligatoire pour la politesse.
\item \eng{Could you wait a moment, please?} — avec un client, registre formel.
\end{enumerate}
\textbf{Conclusion :} la structure est invariable : \fbox{Could/Can you + base verbale (+ complément) (+ please)?}. Le choix \eng{could/can} dépend du rapport social avec l'interlocuteur.
\end{exoresolu}

\courssub{Fiche méthode 4 — Traduire sans calquer le français}

\begin{methode}[Éviter les calques du français dans les demandes polies]
\begin{enumerate}
\item « Est-ce que tu peux... ? » se traduit par \eng{Can you...?}, jamais par \eng{Is it that you can...?}.
\item « Pourriez-vous... ? » = \eng{Could you...?}. Ne traduis pas « pourriez » par un conditionnel complet du type \eng{would you can}.
\item « Laisse-moi t'aider » = \eng{Let me help you.} Le verbe \eng{let} est suivi de la base verbale, contrairement au français « laisser quelqu'un \emph{faire} » qui n'a pas de \eng{to} en anglais non plus — mais attention à ne pas mettre \eng{to} par analogie avec d'autres verbes.
\item « S'il vous plaît » = \eng{please}, placé en fin de demande ou après \eng{you} ; jamais \eng{if it pleases you} dans une demande courante.
\item Vérifie toujours l'ordre : \fbox{Could you + BV} et non \fbox{You could + BV?} (qui reste compréhensible mais sonne comme un calque).
\end{enumerate}
\end{methode}

\begin{exoresolu}[Traduire en anglais correct]
\textbf{Énoncé.} Translate into English.\\
1. Pourriez-vous m'indiquer la banque la plus proche, s'il vous plaît ?\\
2. Laisse-moi payer le taxi.\\
3. Tu peux m'aider à porter ces cartons ?\\
4. Laissons-nous commencer par les questions des clients. (suggestion : commençons)\\
5. Pourriez-vous parler plus lentement, s'il vous plaît ?
\tcblower
\textbf{Solution détaillée.}
\begin{enumerate}
\item \eng{Could you tell me the way to the nearest bank, please?} — « indiquer » se rend naturellement par \eng{tell me the way to} ; \eng{nearest} = « le plus proche » (superlatif).
\item \eng{Let me pay for the taxi.} — attention : \eng{pay \textbf{for} something} ; on peut aussi dire \eng{Let me pay the taxi fare.}
\item \eng{Can you help me (to) carry these boxes?} — registre familier $\Rightarrow$ \eng{can} ; après \eng{help}, le \eng{to} est facultatif.
\item \eng{Let's start with the customers' questions.} — suggestion collective $\Rightarrow$ \eng{let's} ; apostrophe du possessif : \eng{customers'}.
\item \eng{Could you speak more slowly, please?} — adverbe \eng{slowly} (pas d'adjectif seul ici en registre soigné) ; comparatif \eng{more slowly}.
\end{enumerate}
\textbf{Conclusion :} la traduction correcte exige de penser la \emph{structure anglaise} (\eng{let me + BV}, \eng{could you + BV}) et non de transposer mot à mot la phrase française.
\end{exoresolu}

\courssub{Fiche méthode 5 — Réagir à une demande polie : accepter ou refuser}

\begin{methode}[Répondre à une demande polie en situation orale]
\begin{enumerate}
\item \textbf{Accepte} avec une formule courte et chaleureuse : \eng{Sure. / Of course. / Certainly. / With pleasure. / No problem.}
\item \textbf{Refuse} toujours avec une excuse : \eng{I'm sorry, I can't. / I'm afraid I can't. / I'd like to, but...} — un refus sec (\eng{No.}) est impoli en anglais.
\item \textbf{Réponds à une offre d'aide} (\eng{Let me help you.}) par : \eng{Thank you, that's very kind of you.} (accepter) ou \eng{No, thank you, I can manage.} (décliner poliment).
\item En production écrite ou orale d'examen, enchaîne demande + réponse pour montrer la maîtrise du dialogue : \eng{— Could you open an account for me, please? — Certainly, madam. Please fill in this form.}
\end{enumerate}
\end{methode}

\begin{exoresolu}[Dialogue complet à la banque]
\textbf{Énoncé.} Complete the dialogue with polite requests, offers and answers.\\
\textbf{Clerk:} Good morning, sir. (1) \dots\dots\dots\ (offer to help)\\
\textbf{Customer:} Good morning. (2) \dots\dots\dots\ (ask politely to open a savings account)\\
\textbf{Clerk:} Certainly, sir. (3) \dots\dots\dots\ (ask the customer to fill in a form)\\
\textbf{Customer:} Of course. And (4) \dots\dots\dots\ (ask politely about the interest rate)\\
\textbf{Clerk:} (5) \dots\dots\dots\ (offer to explain everything)
\tcblower
\textbf{Solution détaillée.}\\
(1) \eng{Let me help you.} ou \eng{How may I help you?} — l'employé offre son aide $\Rightarrow$ \eng{let me}.\\
(2) \eng{Could you open a savings account for me, please?} — client face à un employé : registre poli $\Rightarrow$ \eng{could you + BV}.\\
(3) \eng{Could you fill in this form, please?} — demande polie formelle ; verbe à particule \eng{fill in}.\\
(4) \eng{could you tell me about the interest rate, please?} — enchaînement après \eng{and} : on garde la minuscule à \eng{could} car la phrase continue.\\
(5) \eng{Let me explain everything to you.} — offre d'aide $\Rightarrow$ \eng{let me + BV} ; préposition \eng{explain something \textbf{to} somebody}.\\
\textbf{Conclusion :} un dialogue d'examen réussi alterne \eng{let me} (offres) et \eng{could you} (demandes), avec des réponses courtes et polies (\eng{certainly}, \eng{of course}). C'est cette alternance qui prouve la maîtrise de la leçon.
\end{exoresolu}

\begin{attention}[Les 5 erreurs qui coûtent des points au Bac]
\begin{enumerate}
\item \textbf{\eng{to} après \eng{let me} :} \eng{Let me to help you.} est faux. Après \eng{let me / let's}, toujours la \emph{base verbale} : \eng{Let me help you.}
\item \textbf{Verbe conjugué après \eng{could you} :} \eng{Could you helps me?} ou \eng{Could you to help me?} sont faux. \eng{could} est un auxiliaire modal : le verbe qui suit reste à la base verbale : \eng{Could you help me?}
\item \textbf{Calque de « est-ce que tu peux » :} ne traduis jamais mot à mot. La demande polie anglaise commence directement par \eng{Can you...?} / \eng{Could you...?}.
\item \textbf{Confusion \eng{lend / borrow} :} \eng{Can you borrow me your pen?} est faux. On \emph{prête} avec \eng{lend somebody something} ; on \emph{emprunte} avec \eng{borrow something from somebody} : \eng{Can you lend me your pen?}
\item \textbf{Refus sec impoli :} répondre \eng{No.} ou \eng{I can't.} seul à une demande polie est sanctionné à l'oral comme à l'écrit. Ajoute toujours une excuse : \eng{I'm sorry, I can't.} / \eng{I'm afraid I can't.}
\end{enumerate}
\end{attention}

\newpage

\coursec{Exercices}

\courssub{Je vérifie mes connaissances}

\begin{exercice}[QCM : la bonne forme]
Pour chaque phrase, choisis la réponse correcte (A, B ou C).
\begin{enumerate}
\item Let me \textbf{...} you with your luggage, madam.\\
A. to help \hspace{1cm} B. help \hspace{1cm} C. helping
\item \textbf{...} you lend me five hundred francs, please?\\
A. Do \hspace{1cm} B. Are \hspace{1cm} C. Could
\item Could you \textbf{...} me the way to the market, please?\\
A. to tell \hspace{1cm} B. telling \hspace{1cm} C. tell
\item \textbf{...} visit the cocoa factory next week!\\
A. Let \hspace{1cm} B. Let's \hspace{1cm} C. Let we
\item Can you \textbf{...} the invoice today, please?\\
A. send \hspace{1cm} B. to send \hspace{1cm} C. sending
\end{enumerate}
\end{exercice}

\begin{exercice}[Vrai ou faux ?]
Indique si chaque affirmation est \textbf{vraie} ou \textbf{fausse}, puis justifie ta réponse en une phrase.
\begin{enumerate}
\item After \eng{let me}, the verb takes \eng{to}.
\item \eng{Could you} is more polite than \eng{can you}.
\item The word \eng{please} can be placed at the beginning or at the end of a request.
\item \eng{Let's} is used to make a suggestion that includes the speaker.
\item In a polite request, the verb after \eng{could you} ends in \eng{-ing}.
\end{enumerate}
\end{exercice}

\begin{exercice}[Vocabulaire : les formules de politesse]
Complète chaque définition avec le mot ou l'expression anglaise qui convient, choisi dans la liste : \eng{request} — \eng{suggestion} — \eng{polite} — \eng{customer} — \eng{invoice}.
\begin{enumerate}
\item A person who buys goods or services in a shop: a \textbf{...}
\item A question you ask when you want someone to do something for you: a \textbf{...}
\item An idea you give about what people could do: a \textbf{...}
\item A document that shows how much money you must pay: an \textbf{...}
\item The opposite of \eng{rude}: \textbf{...}
\end{enumerate}
\end{exercice}

\begin{exercice}[Conjugaison et formes à compléter]
Complète les phrases avec la forme correcte du verbe entre parenthèses (base verbale).
\begin{enumerate}
\item Let me \textbf{...} (carry) your bags to the taxi.
\item Could you \textbf{...} (speak) more slowly, please?
\item Let's \textbf{...} (open) a savings account at the bank.
\item Can you \textbf{...} (show) me your price list, please?
\item Let me \textbf{...} (check) the total for you, sir.
\end{enumerate}
\end{exercice}

\courssub{Je m'entraîne}

\begin{exercice}[Traduction]
Traduis les phrases suivantes en anglais, en utilisant la structure demandée entre parenthèses.
\begin{enumerate}
\item Laissez-moi vous aider, madame. (\eng{let me})
\item Pourriez-vous fermer la porte, s'il vous plaît ? (\eng{could you})
\item Allons au marché de Mokolo cet après-midi ! (\eng{let's})
\item Pouvez-vous me donner votre numéro de téléphone, s'il vous plaît ? (\eng{can you})
\item Laisse-moi payer le taxi-moto. (\eng{let me})
\end{enumerate}
\end{exercice}

\begin{exercice}[Transformation : des ordres aux demandes polies]
Transforme chaque ordre en demande polie avec \eng{could you ... please} ou \eng{can you ... please}.\\
Exemple : \eng{Close the window.} $\rightarrow$ \eng{Could you close the window, please?}
\begin{enumerate}
\item Give me the receipt.
\item Open the shop at eight o'clock.
\item Sign this form.
\item Call the supplier in Douala.
\item Bring me a bag, please.
\end{enumerate}
\end{exercice}

\begin{exercice}[Texte à trous : à la banque]
Complète le dialogue suivant avec les mots de la liste : \eng{let me} — \eng{let's} — \eng{could you} — \eng{can you} — \eng{please} — \eng{let me}.
\end{exercice}

\begin{textebox}[At the bank in Yaoundé]
\textbf{Customer:} Good morning, sir. (1) \textbf{...} open a savings account, \textbf{...} (2)?\\
\textbf{Clerk:} Of course, madam. (3) \textbf{...} see your identity card.\\
\textbf{Customer:} Here it is. And (4) \textbf{...} you explain the interest rate to me, please?\\
\textbf{Clerk:} Certainly. (5) \textbf{...} sit down and look at this leaflet together.\\
\textbf{Customer:} Thank you. (6) \textbf{...} fill in the form now.
\end{textebox}

\begin{exercice}[Appariement : demande et réponse]
Associe chaque demande polie (1--5) à la réponse qui convient (a--e).
\begin{enumerate}
\item Could you lend me your calculator, please?
\item Let me help you carry those boxes.
\item Can you tell me the price of this bag of rice?
\item Let's meet at the market at ten o'clock.
\item Could you repeat the last figure, please?
\end{enumerate}
\begin{itemize}
\item[a.] \eng{Good idea. See you there!}
\item[b.] \eng{Of course. It is fifteen thousand francs.}
\item[c.] \eng{Sure, here you are.}
\item[d.] \eng{Thank you, that is very kind of you.}
\item[e.] \eng{Yes, certainly: two million, five hundred thousand.}
\end{itemize}
\end{exercice}

\courssub{Je me prépare au Bac}

\begin{exercice}[Compréhension et langue : texte]
Lis le texte suivant, puis réponds aux questions en anglais, avec des phrases complètes.
\end{exercice}

\begin{textebox}[A polite shopkeeper in Buea]
Mrs Efosi runs a small grocery shop near the market in Buea. Every morning, she opens her shop at seven o'clock and welcomes her customers with a smile. She believes that politeness is good for business. When a customer enters, she often says: “Let me help you find what you need.” If the shop is busy, she asks kindly: “Could you wait a moment, please?”

Her son, who is a student, sometimes works in the shop during the holidays. Mrs Efosi teaches him to speak politely to everyone. “Never shout at a customer,” she tells him. “Say ‘please’ and ‘thank you’, and people will come back.” Thanks to this attitude, her little shop is always full, and many customers have become her friends.
\end{textebox}

\textbf{Questions :}
\begin{enumerate}
\item Where is Mrs Efosi's shop? (1 mark)
\item What time does she open her shop every morning? (1 mark)
\item Why does Mrs Efosi think politeness is important? (2 marks)
\item Find in the text one polite request with \eng{could you} and copy it. (2 marks)
\item Find in the text one offer of help with \eng{let me} and copy it. (2 marks)
\item Transform into a polite request: \eng{Wait a moment!} (2 marks)
\end{enumerate}

\begin{exercice}[Traduction type Bac]
Traduis le passage suivant en anglais.\\
« Bonjour, monsieur. Pourriez-vous m'aider, s'il vous plaît ? Je voudrais envoyer de l'argent à mon frère à Bamenda. Laissez-moi remplir ce formulaire. Pouvez-vous me dire combien coûte le transfert ? Merci beaucoup pour votre aide. »
\end{exercice}

\begin{exercice}[Production écrite type Bac]
You are the manager of a small bookshop in Douala. Write a polite e-mail (about 80--100 words) to your supplier, Mr Johnson, in which you:
\begin{itemize}
\item greet him politely;
\item make at least \textbf{three polite requests} using \eng{could you}, \eng{can you} or \eng{please} (ask for a price list, a delivery date, a catalogue...);
\item make \textbf{one suggestion} with \eng{let's} (for example, meeting next week);
\item end the e-mail politely.
\end{itemize}
\end{exercice}

\newpage

\coursec{Corrigés des exercices}

\begin{corrige}[Exercice 1 — QCM : la bonne forme]
\begin{enumerate}
\item \textbf{B. help} — \eng{Let me help you with your luggage, Madam.} Après \eng{let me}, on utilise toujours la \cle{base verbale} (infinitif sans \eng{to}). Jamais \eng{to help} ni \eng{helping}. \textit{Note : \eng{Madam} (adresse directe) prend une majuscule.}
\item \textbf{C. Could} — \eng{Could you lend me five hundred francs, please?} C'est la formule de demande polie ; \eng{Do you lend} signifierait « Prêtes-tu (habituellement) ? » et \eng{Are you lend} est impossible.
\item \textbf{C. tell} — \eng{Could you tell me the way to the market, please?} Après \eng{could you}, le verbe reste à la base verbale, comme après tout modal.
\item \textbf{B. Let's} — \eng{Let's visit the cocoa factory next week!} \eng{Let's} = \eng{let us} (suggestion incluant le locuteur). \eng{Let we} n'existe pas : après \eng{let}, on met le pronom \emph{complément} (\eng{us}, pas \eng{we}).
\item \textbf{A. send} — \eng{Can you send the invoice today, please?} Même règle : modal \eng{can} + base verbale.
\end{enumerate}
\end{corrige}

\begin{corrige}[Exercice 2 — Vrai ou faux ?]
\begin{enumerate}
\item \textbf{Faux.} After \eng{let me}, the verb is in the bare infinitive, without \eng{to}: \eng{Let me help you} (et non \eng{let me to help}).
\item \textbf{Vrai.} \eng{Could you} is more polite and more formal than \eng{can you}; \eng{can you} is used with friends or in informal situations.
\item \textbf{Vrai.} \eng{Please} can come at the beginning or at the end: \eng{Please open the door.} / \eng{Open the door, please.} (avec une virgule en fin de phrase).
\item \textbf{Vrai.} \eng{Let's} (= \eng{let us}) is used for a suggestion that includes the speaker: \eng{Let's go to the market!} (« Allons au marché ! »).
\item \textbf{Faux.} After \eng{could you}, the verb is in the bare infinitive: \eng{Could you repeat, please?} — jamais de forme en \eng{-ing} après un modal.
\end{enumerate}
\end{corrige}

\begin{corrige}[Exercice 3 — Vocabulaire : les formules de politesse]
\begin{enumerate}
\item a \textbf{customer} — « un client » (attention au faux ami : \eng{customer} = client d'un magasin ; \eng{costumer} n'existe pas dans ce sens, c'est un \emph{costumier}).
\item a \textbf{request} — « une demande » (nom ; le verbe est \eng{to request}, registre soutenu).
\item a \textbf{suggestion} — « une suggestion » ; le verbe correspondant est \eng{to suggest}.
\item an \textbf{invoice} — « une facture » (on dit \eng{an invoice} car le mot commence par une voyelle).
\item \textbf{polite} — « poli » ; contraires : \eng{rude} ou \eng{impolite}.
\end{enumerate}
\end{corrige}

\begin{corrige}[Exercice 4 — Conjugaison et formes à compléter]
Dans tous les cas, la structure exige la \cle{base verbale} : après \eng{let me}, \eng{let's}, \eng{could you} et \eng{can you}, le verbe ne change jamais (pas de \eng{to}, pas de \eng{-s}, pas de \eng{-ing}).
\begin{enumerate}
\item Let me \textbf{carry} your bags to the taxi. « Laisse-moi porter tes sacs jusqu'au taxi. »
\item Could you \textbf{speak} more slowly, please? « Pourriez-vous parler plus lentement, s'il vous plaît ? »
\item Let's \textbf{open} a savings account at the bank. « Ouvrons un compte d'épargne à la banque. »
\item Can you \textbf{show} me your price list, please? « Pouvez-vous me montrer votre liste de prix, s'il vous plaît ? »
\item Let me \textbf{check} the total for you, Sir. « Laissez-moi vérifier le total pour vous, Monsieur. » \textit{Note : \eng{Sir} (adresse directe) prend une majuscule.}
\end{enumerate}
\end{corrige}

\begin{corrige}[Exercice 5 — Traduction]
\begin{enumerate}
\item \eng{Let me help you, Madam.} — Pas de \eng{to} après \eng{let me} ; \eng{Madam} (vocatif) prend une majuscule.
\item \eng{Could you close the door, please?} — \eng{Could you} rend le « pourriez-vous » du français ; virgule avant \eng{please} en fin de phrase.
\item \eng{Let's go to the Mokolo market this afternoon!} — \eng{Let's} traduit « allons » ; ne pas oublier le verbe \eng{go} : on ne dit pas \eng{let's to the market}.
\item \eng{Can you give me your phone number, please?} — Ordre des mots : \eng{give me your phone number} (donner quelque chose à quelqu'un = \eng{give somebody something}).
\item \eng{Let me pay for the motorbike taxi.} — Attention : \eng{pay for something} (payer quelque chose) ; \eng{pay the taxi} signifierait payer le chauffeur lui-même.
\end{enumerate}
\end{corrige}

\begin{corrige}[Exercice 6 — Transformation : des ordres aux demandes polies]
\begin{enumerate}
\item \eng{Could you give me the receipt, please?}
\item \eng{Can you open the shop at eight o'clock, please?} (ou \eng{Could you open...})
\item \eng{Could you sign this form, please?}
\item \eng{Could you call the supplier in Douala, please?}
\item \eng{Can you bring me a bag, please?}
\end{enumerate}
\textbf{Méthode :} on supprime l'impératif, on place \eng{Could you} ou \eng{Can you} en tête, on garde le verbe à la base verbale, et on ajoute \eng{please} en fin de phrase précédé d'une virgule. La phrase devient interrogative : point d'interrogation obligatoire.
\end{corrige}

\begin{corrige}[Exercice 7 — Texte à trous : à la banque]
\begin{enumerate}
\item \textbf{Can you} — \eng{Can you open a savings account...} : demande polie adressée à l'employé.
\item \textbf{please} — placé en fin de demande, après la virgule.
\item \textbf{Let me} — \eng{Let me see your identity card.} : offre/demande d'autorisation de l'employé (« Laissez-moi voir... »).
\item \textbf{could} — \eng{And could you explain the interest rate to me, please?} : demande plus formelle du client (\eng{explain sth to sb}).
\item \textbf{Let's} — \eng{Let's sit down and look at this leaflet together.} : suggestion incluant les deux personnes, confirmée par \eng{together}.
\item \textbf{Let me} — \eng{Let me fill in the form now.} : le client propose de faire l'action lui-même (\eng{fill in} = remplir, UK).
\end{enumerate}
\textbf{Dialogue complet :}\\
\eng{“Good morning, Sir. Can you open a savings account, please?”} \\
\eng{“Of course, Madam. Let me see your identity card.”} \\
\eng{“Here it is. And could you explain the interest rate to me, please?”} \\
\eng{“Certainly. Let's sit down and look at this leaflet together.”} \\
\eng{“Thank you. Let me fill in the form now.”}
\end{corrige}

\begin{corrige}[Exercice 8 — Appariement : demande et réponse]
\begin{center}
\begin{tabular}{|l|l|}
\hline
\rowcolor{popblueL} \textbf{Demande} & \textbf{Réponse} \\
\hline
1. Could you lend me your calculator, please? & \textbf{c.} Sure, here you are. \\
\hline
2. Let me help you carry those boxes. & \textbf{d.} Thank you, that is very kind of you. \\
\hline
3. Can you tell me the price of this bag of rice? & \textbf{b.} Of course. It is fifteen thousand francs. \\
\hline
4. Let's meet at the market at ten o'clock. & \textbf{a.} Good idea. See you there! \\
\hline
5. Could you repeat the last figure, please? & \textbf{e.} Yes, certainly: two million, five hundred thousand. \\
\hline
\end{tabular}
\end{center}
\textbf{Justifications :} (1) on prête un objet → \eng{here you are} « voilà » ; (2) on remercie pour une offre d'aide ; (3) la réponse donne un prix ; (4) on accepte une suggestion → \eng{Good idea} ; (5) la réponse répète un chiffre (\eng{figure} = chiffre, nombre).
\end{corrige}

\begin{corrige}[Exercice 9 — Compréhension et langue : texte]
\textbf{Barème indicatif : 10 points.}
\begin{enumerate}
\item \eng{Her shop is near the market in Buea.} (1 point — réponse complète avec \eng{near the market} et \eng{Buea}.)
\item \eng{She opens her shop at seven o'clock every morning.} (1 point — attention à la 3\textsuperscript{e} personne : \eng{opens}.)
\item \eng{She thinks politeness is important because it is good for business: when you say “please” and “thank you”, customers come back.} (2 points — l'idée : la politesse fidélise la clientèle.)
\item \eng{“Could you wait a moment, please?”} (2 points — recopier exactement la demande du texte.)
\item \eng{“Let me help you find what you need.”} (2 points — recopier exactement l'offre d'aide.)
\item \eng{Could you wait a moment, please?} ou \eng{Can you wait a moment, please?} (2 points — transformation de l'impératif en demande polie.)
\end{enumerate}
\textbf{Points de vigilance :} répondre par des phrases complètes avec sujet et verbe conjugué ; ne pas recopier tout le texte ; respecter la base verbale après \eng{could you} à la question 6.
\end{corrige}

\begin{corrige}[Exercice 10 — Traduction type Bac]
\textbf{Traduction modèle :}\\
\eng{“Good morning, Sir. Could you help me, please? I would like to send some money to my brother in Bamenda. Let me fill in this form. Can you tell me how much the transfer costs? Thank you very much for your help.”}

\textbf{Barème indicatif (10 points) :} 2 points par phrase correctement traduite (5 phrases).
\begin{itemize}
\item « Pourriez-vous » → \eng{Could you} (registre poli soutenu) ; « Pouvez-vous » → \eng{Can you}.
\item « Je voudrais » → \eng{I would like to} + base verbale (et non \eng{I want}, trop direct).
\item « envoyer de l'argent à mon frère » → \eng{send some money to my brother} (préposition \eng{to}).
\item « remplir ce formulaire » → \eng{fill in this form} (verbe à particule ; \eng{fill out} est accepté, américain).
\item « combien coûte le transfert » → \eng{how much the transfer costs} : attention à l'\cle{interrogation indirecte} — pas d'inversion, pas de \eng{does} : on ne dit pas \eng{how much does the transfer cost} après \eng{Can you tell me}.
\end{itemize}
\end{corrige}

\begin{corrige}[Exercice 11 — Production écrite type Bac]
\textbf{Barème indicatif (10 points) :} respect du format e-mail et de la longueur (2 pts) ; au moins trois demandes polies correctes (3 pts) ; une suggestion avec \eng{let's} (2 pts) ; ouverture et clôture polies (1 pt) ; correction grammaticale et vocabulaire (2 pts).

\textbf{Proposition de rédaction modèle :}\par
\textbf{Subject:} Book order — price list and delivery\\[0.3em]
\eng{Dear Mr Johnson,}\\
\eng{I hope you are well. I am writing about our next book order. Could you send me your new price list, please? Can you also give me the delivery date for the novels we ordered last month? And could you send us your latest catalogue, please?}\\
\eng{Let's meet next week to discuss the order. Please let me know which day suits you best.}\\
\eng{Thank you very much for your help.}\\
\eng{Best regards,}\\
\eng{Amina Mbappe}\\
\eng{Manager, Rainbow Bookshop, Douala}

\textbf{Points de vigilance :}
\begin{itemize}
\item Base verbale après \eng{could you} / \eng{can you} : \eng{Could you send...} (jamais \eng{to send}).
\item \eng{Please} en fin de demande, précédé d'une virgule.
\item \eng{Let's} + base verbale pour la suggestion : \eng{Let's meet next week.}
\item Formules de clôture : \eng{Best regards} ou \eng{Yours sincerely} ; éviter \eng{Bye} dans un e-mail professionnel.
\item Compter les mots (80--100) : le modèle ci-dessus en contient environ 90.
\end{itemize}
\end{corrige}

\newpage

\coursec{Fiche bilan}

\begin{popfigure}
\begin{center}\tcbox[colback=popblue,colframe=popblue,coltext=white,arc=9pt,boxsep=4pt,left=6pt,right=6pt]{\bfseries\small Polite requests \& suggestions}\end{center}
\vspace{-2pt}
\begin{tcbraster}[raster columns=3,raster equal height=rows,raster column skip=6pt,raster row skip=6pt,enhanced,blankest]
\begin{tcolorbox}[enhanced,colback=white,colframe=poporange,boxrule=0.9pt,arc=3pt,title={Let me + BV},fonttitle=\bfseries\scriptsize,coltitle=white,colbacktitle=poporange,left=4pt,right=4pt,top=2pt,bottom=2pt,boxsep=1pt,toptitle=1.5pt,bottomtitle=1.5pt,halign=left]
\begin{itemize}[nosep,leftmargin=*,itemsep=1pt,font=\scriptsize]
\item offre d'aide
\item suggestion
\end{itemize}
\end{tcolorbox}
\begin{tcolorbox}[enhanced,colback=white,colframe=popgreen,boxrule=0.9pt,arc=3pt,title={Could you + BV ?},fonttitle=\bfseries\scriptsize,coltitle=white,colbacktitle=popgreen,left=4pt,right=4pt,top=2pt,bottom=2pt,boxsep=1pt,toptitle=1.5pt,bottomtitle=1.5pt,halign=left]
\begin{itemize}[nosep,leftmargin=*,itemsep=1pt,font=\scriptsize]
\item très poli
\end{itemize}
\end{tcolorbox}
\begin{tcolorbox}[enhanced,colback=white,colframe=poppurple,boxrule=0.9pt,arc=3pt,title={Can you + BV ?},fonttitle=\bfseries\scriptsize,coltitle=white,colbacktitle=poppurple,left=4pt,right=4pt,top=2pt,bottom=2pt,boxsep=1pt,toptitle=1.5pt,bottomtitle=1.5pt,halign=left]
\begin{itemize}[nosep,leftmargin=*,itemsep=1pt,font=\scriptsize]
\item courant
\end{itemize}
\end{tcolorbox}
\begin{tcolorbox}[enhanced,colback=white,colframe=poppink,boxrule=0.9pt,arc=3pt,title={Please},fonttitle=\bfseries\scriptsize,coltitle=white,colbacktitle=poppink,left=4pt,right=4pt,top=2pt,bottom=2pt,boxsep=1pt,toptitle=1.5pt,bottomtitle=1.5pt,halign=left]
\begin{itemize}[nosep,leftmargin=*,itemsep=1pt,font=\scriptsize]
\item début ou fin
\end{itemize}
\end{tcolorbox}
\begin{tcolorbox}[enhanced,colback=white,colframe=popgold,boxrule=0.9pt,arc=3pt,title={Réponses},fonttitle=\bfseries\scriptsize,coltitle=white,colbacktitle=popgold,left=4pt,right=4pt,top=2pt,bottom=2pt,boxsep=1pt,toptitle=1.5pt,bottomtitle=1.5pt,halign=left]
\begin{itemize}[nosep,leftmargin=*,itemsep=1pt,font=\scriptsize]
\item Sure / Of course
\item I'm sorry, but...
\end{itemize}
\end{tcolorbox}
\begin{tcolorbox}[enhanced,colback=white,colframe=popblue!60,boxrule=0.9pt,arc=3pt,title={Pièges},fonttitle=\bfseries\scriptsize,coltitle=white,colbacktitle=popblue!60,left=4pt,right=4pt,top=2pt,bottom=2pt,boxsep=1pt,toptitle=1.5pt,bottomtitle=1.5pt,halign=left]
\begin{itemize}[nosep,leftmargin=*,itemsep=1pt,font=\scriptsize]
\item pas de to
\item pas de -s
\end{itemize}
\end{tcolorbox}
\end{tcbraster}

\legende{Carte mentale de la leçon : demandes et suggestions polies}
\end{popfigure}

\begin{bilan}
\textbf{1. Lexique clé (contexte économique)}

\begin{center}
\begin{tabularx}{\linewidth}{|X|X|}\hline
\rowcolor{popblueL} \textbf{English} & \textbf{Français} \\ \hline
a request & une demande \\ \hline
a suggestion & une suggestion \\ \hline
to lend (lent, lent) & prêter (à qqn) \\ \hline
to borrow & emprunter (à qqn) \\ \hline
to help somebody & aider quelqu'un \\ \hline
to carry the goods & porter les marchandises \\ \hline
to give somebody a discount & faire une remise à quelqu'un \\ \hline
polite / politely & poli / poliment \\ \hline
\end{tabularx}
\end{center}

\textbf{2. Règles de grammaire à connaître par cœur}

\begin{center}
\begin{tabularx}{\linewidth}{|l|X|X|}\hline
\rowcolor{popblueL} \textbf{Structure} & \textbf{Emploi} & \textbf{Exemple} \\ \hline
\cle{Let me} + base verbale & offrir son aide, suggérer & \eng{Let me help you.} \\ \hline
\cle{Can you} + BV ? & demande polie courante & \eng{Can you open the door?} \\ \hline
\cle{Could you} + BV ? & demande plus polie, formelle & \eng{Could you lend me 500 francs?} \\ \hline
\cle{Please} (début ou fin) & adoucit toute demande & \eng{Pass me the salt, please.} \\ \hline
\end{tabularx}
\end{center}

\begin{itemize}
\item Après \eng{let me}, \eng{can you} et \eng{could you}, le verbe est TOUJOURS à la \cle{base verbale} : jamais de \eng{to}, jamais de \eng{-s}, jamais de \eng{-ing}.
\item \eng{Could you} n'est PAS un passé ici : c'est simplement plus poli que \eng{can you}.
\item Réponses typiques : \eng{Sure. / Of course. / Certainly.} (oui) ; \eng{I'm sorry, but I can't.} (non poli).
\item \eng{Please} se place au début ou à la fin de la phrase, jamais entre le sujet et le verbe.
\end{itemize}

\textbf{3. Méthode express}
\begin{enumerate}
\item Identifier l'intention : offrir de l'aide (\eng{let me}) ou demander un service (\eng{can/could you}).
\item Choisir le degré de politesse : ami $\rightarrow$ \eng{can you} ; inconnu, client, supérieur $\rightarrow$ \eng{could you} + \eng{please}.
\item Construire : structure + base verbale, sans \eng{to}.
\item Vérifier la réponse : accepter (\eng{Sure!}) ou refuser poliment (\eng{I'm sorry, but...}).
\end{enumerate}
\end{bilan}

\begin{aretenir}[Checklist avant le Bac]
\begin{itemize}
\item[$\square$] Je sais conjuguer la structure \eng{let me} + base verbale sans \eng{to}.
\item[$\square$] Je distingue \eng{can you} (courant) et \eng{could you} (plus poli).
\item[$\square$] Je n'écris jamais \eng{Could you to help me?} : pas de \eng{to} après \eng{could}.
\item[$\square$] Je n'ajoute jamais de \eng{-s} au verbe : \eng{Can you help} (et non \eng{helps}).
\item[$\square$] Je place \eng{please} au début ou à la fin de la demande.
\item[$\square$] Je sais répondre poliment : \eng{Sure. / Of course. / I'm sorry, but I can't.}
\item[$\square$] Je ne traduis pas mot à mot « pouvez-vous » par \eng{are you able to} dans une demande simple.
\item[$\square$] Je sais transformer un ordre direct (\eng{Help me!}) en demande polie (\eng{Could you help me, please?}).
\item[$\square$] Je sais utiliser ces structures dans un dialogue au marché, à la banque ou au bureau.
\item[$\square$] Je me souviens que \eng{could} n'exprime pas le passé dans une demande polie.
\end{itemize}
\end{aretenir}

\findecours

\end{document}
